% Pengaturan ukuran teks dan bentuk halaman dua sisi
\documentclass[12pt]{book}

% Pengaturan ukuran halaman dan margin
\usepackage[a4paper,top=30mm,left=30mm,right=20mm,bottom=25mm]{geometry}

% Pengaturan bahasa Indonesia untuk pemotongan kata (Hyphenation) & Rata Kanan-Kiri
\usepackage[indonesian]{babel}
\usepackage{ragged2e}

% Pengaturan ukuran spasi
\usepackage[singlespacing]{setspace}

% Pengaturan caption untuk tabel
\usepackage{caption}

% Judul dokumen
\title{Proposal Tugas Akhir ITS}
\author{Rakan, Sultan Syafiq}

% Pengaturan detail pada file PDF
\usepackage[pdfauthor={Rakan, Sultan Syafiq},bookmarksnumbered,pdfborder={0 0 0}]{hyperref}

% Pengaturan ukuran indentasi
% Paragraf rata kanan-kiri tanpa inden baris pertama; antarparagraf diberi jarak
\usepackage[skip=6pt plus 2pt, indent=0pt]{parskip}
\setlength{\JustifyingParindent}{0pt} % \justifying (ragged2e) memakai nilai ini, bukan \parindent
\usepackage{microtype} % Merapikan spasi antarkata agar tidak ada baris yang keluar margin kanan
\setlength{\emergencystretch}{1.5em}

% Package lainnya
\usepackage{changepage}
\usepackage{etoolbox} % Mengubah fungsi default

% Pengaturan jenis karakter & Font Times Modern
\usepackage[utf8]{inputenc}
\usepackage{amsmath} % Untuk environment cases dan persamaan
\usepackage{newtxtext,newtxmath} % Pengganti txfonts yang modern & stabil

\usepackage[style=ieee, backend=biber, safeinputenc]{biblatex} % string bahasa dari indonesian.lbx
\usepackage{enumitem} % Pembuatan list
\usepackage{lipsum} % Pembuatan template kalimat
\usepackage{graphicx} % Input gambar
\usepackage{longtable} % Pembuatan tabel
\usepackage[table,xcdraw]{xcolor} % Pewarnaan tabel
\usepackage{eso-pic} % Untuk menggunakan background image di halaman
\usepackage{multicol} % Pembuatan kolom ganda
\usepackage{multirow} % Pembuatan baris ganda
\usepackage{tabularx} % Untuk mengatur kolom, seperti grid pada CSS
\usepackage{wrapfig}
\usepackage{float}
\usepackage{tikz} % Flowchart penelitian
\usetikzlibrary{positioning,shapes.geometric,arrows.meta}

% Pengaturan format daftar isi, daftar gambar, dan daftar tabel
\usepackage[titles]{tocloft}
\setlength{\cftsecindent}{2em}
\setlength{\cftsubsecindent}{2em}
\setlength{\cftbeforechapskip}{1.5ex}
\setlength{\cftbeforesecskip}{1.5ex}
\setlength{\cftbeforetoctitleskip}{0cm}
\setlength{\cftbeforeloftitleskip}{0cm}
\setlength{\cftbeforelottitleskip}{0cm}
\renewcommand{\cfttoctitlefont}{\hfill\Large\bfseries} % command untuk membuat heading bold dan besar
\renewcommand{\cftaftertoctitle}{\hfill}
\renewcommand{\cftloftitlefont}{\hfill\Large\bfseries}
\renewcommand{\cftafterloftitle}{\hfill}
\renewcommand{\cftlottitlefont}{\hfill\Large\bfseries}
\renewcommand{\cftafterlottitle}{\hfill}

% Definisi untuk "Halaman ini sengaja dikosongkan"
\patchcmd{\cleardoublepage}{\hbox{}}{
  \thispagestyle{empty}
  \vspace*{\fill}
  \begin{center}\textit{[Halaman ini sengaja dikosongkan]}\end{center}
  \vfill}{}{}

% Pengaturan penomoran halaman
\usepackage{fancyhdr}
\fancyhf{}
\renewcommand{\headrulewidth}{0pt}
\pagestyle{fancy}
\fancyfoot[C,CO]{\thepage}
\patchcmd{\chapter}{plain}{fancy}{}{}
\patchcmd{\chapter}{empty}{plain}{}{}

% Pengaturan format judul bab
\usepackage{titlesec}
\renewcommand{\thesection}{\thechapter.\arabic{section}}
\titleformat{\chapter}[hang]{\centering\bfseries\large}{BAB\ \arabic{chapter}\ }{0ex}{\vspace{0ex}\centering}
\titleformat*{\section}{\large\bfseries}
\titleformat*{\subsection}{\normalsize\bfseries}
\titlespacing{\chapter}{0ex}{0ex}{4ex}
\titlespacing{\section}{0ex}{1ex}{0ex}
\titlespacing{\subsection}{0ex}{0.5ex}{0ex}
\titleformat{\subsubsection}{\normalsize\bfseries}{\thesubsubsection}{1em}{}
\titlespacing{\section}{0pt}{2ex plus .5ex minus .2ex}{1ex}
\titlespacing{\subsection}{0pt}{1.5ex plus .5ex minus .2ex}{0.8ex}
\titlespacing{\subsubsection}{0pt}{1.2ex plus .4ex minus .2ex}{0.6ex}
\raggedbottom % Matikan peregangan vertikal halaman
\setcounter{secnumdepth}{3} % Untuk memberi penomoran pada \subsubsection

\counterwithin{figure}{chapter}
\counterwithin{table}{chapter}

% Mengganti figure dan table menjadi gambar dan tabel
\renewcommand{\figurename}{Gambar}
\renewcommand{\tablename}{Tabel}

% !TEX root = ../main.tex
% Tambahkan format tanda hubung yang benar di sini
\hyphenation{
  ro-ket
  me-ngem-bang-kan
  per-hi-tu-ngan
}

% Menambahkan resource daftar pustaka
\addbibresource{pustaka/pustaka.bib}

% Isi keseluruhan dokumen
\begin{document}
  \justifying % Memaksa seluruh dokumen rata kanan-kiri (Justified)

  % Nomor halaman pembuka dimulai dari sini
  \pagenumbering{roman}

  % Atur ulang penomoran halaman
  \setcounter{page}{1}

  % Sampul Bahasa Indonesia
  \newcommand\covercontents{sampul/konten-id.tex}
  % !TEX root = ../main.tex
\AddToShipoutPictureBG*{
  \AtPageLowerLeft{
    % Ubah nilai berikut jika posisi horizontal background tidak sesuai
    \hspace{-3.25mm}

    % Ubah nilai berikut jika posisi vertikal background tidak sesuai
    \raisebox{0mm}{
      \includegraphics[width=\paperwidth,height=\paperheight]{sampul/gambar/sampul-luar-tipis.png}
    }
  }
}

% Menyembunyikan nomor halaman
\thispagestyle{empty}

% Pengaturan margin untuk menyesuaikan konten sampul
\newgeometry{
  top=65mm,
  left=30mm,
  right=30mm,
  bottom=20mm
}

\begin{flushleft}

  % Pemilihan font sans serif
  \sffamily

  % Pemilihan font bold
  \fontseries{bx}
  \selectfont
  \begin{spacing}{1.5}
    \input{\covercontents}
  \end{spacing}

\end{flushleft}

\restoregeometry


  % Lembar pengesahan
  % !TEX root = ../main.tex
\chapter*{LEMBAR PENGESAHAN}

% Menyembunyikan nomor halaman
\thispagestyle{empty}

\begin{center}
  % Ubah kalimat berikut dengan judul tugas akhir
  \textbf{KLASIFIKASI PERILAKU PENGEMUDI AGRESIF PADA KENDARAAN BERBASIS TELEMETRI OBD-II MENGGUNAKAN CNN-BILSTM DAN PRELU}
\end{center}

\begingroup
% Pemilihan font ukuran small
\small

\begin{center}
  % Ubah kalimat berikut dengan pernyataan untuk lembar pengesahan
  \textbf{PROPOSAL TUGAS AKHIR} \\
  Diajukan untuk memenuhi salah satu syarat memperoleh gelar
  Sarjana Teknik pada
  Program Studi S-1 Teknik Komputer \\
  Departemen Teknik Komputer \\
  Fakultas Teknologi Elektro dan Informatika Cerdas \\
  Institut Teknologi Sepuluh Nopember
\end{center}

\begin{center}
  % Ubah kalimat berikut dengan nama dan NRP mahasiswa
  Oleh: \textbf{Sultan Syafiq Rakan} \\
  NRP. 5024231009
\end{center}

\begin{center}
  Disetujui Oleh:
\end{center}

\vspace{10ex}

\begingroup
% Menghilangkan padding
\setlength{\tabcolsep}{0pt}

\noindent
\begin{tabularx}{\textwidth}{X c}
  % Ubah kalimat-kalimat berikut dengan nama dan NIP dosen pembimbing pertama
  Dr. Ahmad Zaini, S.T., M.Sc.      &                 \\
  NIP: 19750419200212 1 003    & (Pembimbing)    \\
\end{tabularx}
\endgroup

\vspace{\fill}

\begin{center}
  Mengetahui,\\
  % Ubah kalimat berikut dengan nama departemen
  Kepala Departemen Teknik Komputer FTEIC - ITS\\
  \vspace{10ex}
  % Ubah kalimat berikut dengan jabatan kepala departemen
  \underline{Dr. Arief Kurniawan, S.T., M.T. }\\
  NIP 18560710 194301 1 001\\
  \vspace{10ex}
  % Ubah text dibawah menjadi tempat dan tanggal
  \textbf{SURABAYA} \\
  \textbf{Mei, 2026}
\end{center}
\endgroup

  \cleardoublepage

  % Abstrak
  % !TEX root = ../main.tex
\chapter*{ABSTRAK}
\begin{center}
  \large
  \textbf{KLASIFIKASI PERILAKU PENGEMUDI AGRESIF PADA KENDARAAN BERBASIS TELEMETRI OBD-II MENGGUNAKAN CNN-BILSTM DAN PRELU}
\end{center}
\addcontentsline{toc}{chapter}{ABSTRAK}
% Menyembunyikan nomor halaman
\thispagestyle{empty}

\begin{flushleft}
  \setlength{\tabcolsep}{0pt}
  \bfseries
  \begin{tabular}{ll@{\hspace{6pt}}l}
  Nama Mahasiswa / NRP&:& Sultan Syafiq Rakan / 5024231009\\
  Departemen&:& Teknik Komputer FTEIC - ITS\\
  Dosen Pembimbing&:& Dr. Ahmad Zaini, S.T., M.Sc.\\
  \end{tabular}
  \vspace{4ex}
\end{flushleft}
\textbf{Abstrak}

% Isi Abstrak
Perilaku mengemudi ugal-ugalan menjadi salah satu penyebab utama kecelakaan lalu lintas, sehingga dibutuhkan sistem yang mengenali perilaku tersebut saat terjadi. Deteksi berbasis kamera membutuhkan komputasi besar dan hanya dapat memantau ruas jalan yang terpasang kamera, sedangkan pemantauan berbasis data kendaraan masih memakai aturan dan ambang batas yang ditentukan secara manual. Penelitian ini mengusulkan klasifikasi perilaku pengemudi dari telemetri kendaraan menggunakan arsitektur hibrida \textit{Convolutional Neural Network} dan \textit{Bidirectional Long Short-Term Memory} (CNN-BiLSTM) dengan fungsi aktivasi \textit{Parametric Rectified Linear Unit} (PReLU). Data latih diambil dari simulator BeamNG.drive melalui protokol UDP OutGauge, dari sesi mengemudi normal dan ugal-ugalan oleh beberapa pengemudi, karena merekam perilaku ugal-ugalan di jalan raya berbahaya. Model kemudian diuji pada mobil nyata yang datanya dibaca melalui port OBD-II dengan adaptor ELM327 WiFi di area tertutup. Agar dapat dipakai pada kedua sumber data, model memakai empat fitur inti yang tersedia di simulator dan di mobil, yaitu kecepatan, putaran mesin, posisi \textit{throttle}, dan akselerasi longitudinal yang dihitung dari selisih kecepatan. Laju cuplik dan resolusi kecepatan data simulator disamakan dengan data OBD-II, lalu data dinormalisasi dengan \textit{Z-score} dan dipotong menjadi jendela tiga detik. Data latih, validasi, dan uji dibagi per blok agar jendela yang tumpang tindih tidak bocor ke data uji. Lapisan konvolusi mengekstrak pola sesaat, lalu lapisan BiLSTM dan LSTM merangkum urutan pola dalam jendela. PReLU mempelajari kemiringan sisi negatifnya selama pelatihan sehingga diharapkan mencegah \textit{dying ReLU}. Model dievaluasi dengan akurasi, presisi, \textit{recall}, \textit{F1-score}, AUC, latensi inferensi, dan waktu deteksi, baik di simulator maupun di mobil. Percobaan pendahuluan di simulator dengan sebelas fitur, satu sesi per kelas, dan satu pengemudi menghasilkan akurasi 97,2\% dan \textit{recall} kelas ugal-ugalan 1,00 pada data uji. Angka ini belum mengukur kemampuan generalisasi, karena data latih dan data uji masih berasal dari sesi yang sama. Penelitian ini diharapkan menjadi dasar sistem pemantauan perilaku pengemudi berbasis OBD-II yang bekerja secara \textit{real-time}.

\vspace{2ex}
\noindent
\textbf{Kata Kunci: \emph{CNN-BiLSTM, PReLU, OBD-II, BeamNG.drive, Perilaku Mengemudi Agresif}}

  \cleardoublepage

  % !TEX root = ../main.tex
\chapter*{ABSTRACT}
\begin{center}
  \large
  \textbf{CLASSIFICATION OF AGGRESSIVE DRIVER BEHAVIOR IN VEHICLES BASED ON OBD-II TELEMETRY USING CNN-BILSTM AND PRELU}
\end{center}
\addcontentsline{toc}{chapter}{ABSTRACT}
% Menyembunyikan nomor halaman
\thispagestyle{empty}

\begin{flushleft}
  \setlength{\tabcolsep}{0pt}
  \bfseries
  \begin{tabular}{lc@{\hspace{6pt}}l}
  Student Name / NRP&: &Sultan Syafiq Rakan / 5024231009\\
  Department&: &Computer Engineering FTEIC - ITS\\
  Advisor&: &Dr. Ahmad Zaini, S.T., M.Sc.\\
  \end{tabular}
  \vspace{4ex}
\end{flushleft}
\textbf{Abstract}

% Isi Abstrak
Aggressive or reckless driving is one of the main causes of traffic accidents, so a system that recognizes this behavior as it happens is needed. Camera-based detection carries a heavy computational load and can only monitor road sections where cameras are installed, while monitoring based on vehicle data still relies on manually defined rules and thresholds. This research proposes driver behavior classification from vehicle telemetry using a hybrid Convolutional Neural Network and Bidirectional Long Short-Term Memory (CNN-BiLSTM) architecture with the Parametric Rectified Linear Unit (PReLU) activation function. Training data are collected from the BeamNG.drive simulator through the OutGauge UDP protocol, from normal and reckless driving sessions by several drivers, because recording reckless driving on public roads is dangerous. The model is then tested on a real car whose data are read from the OBD-II port with an ELM327 WiFi adapter in a closed area. To work on both data sources, the model uses four core features available in the simulator and in the car: speed, engine speed, throttle position, and longitudinal acceleration computed from the change in speed. The sampling rate and speed resolution of the simulator data are matched to the OBD-II data, and the data are then normalized with Z-score and cut into three-second windows. The training, validation, and test sets are split per block so that overlapping windows do not leak into the test set. The convolutional layer extracts short-term patterns, and the BiLSTM and LSTM layers then summarize the sequence of patterns within a window. PReLU learns its negative slope during training, which is expected to prevent dying ReLU. The model is evaluated using accuracy, precision, recall, F1-score, AUC, inference latency, and detection time, both in the simulator and in the car. A preliminary simulator experiment with eleven features, one session per class, and one driver achieved 97.2\% accuracy and a reckless-class recall of 1.00 on the test set. This result does not yet measure generalization, because the training and test data still come from the same session. This research is expected to provide a basis for a real-time driver behavior monitoring system based on OBD-II.

\vspace{2ex}
\noindent
\textbf{Keywords: \emph{CNN-BiLSTM, PReLU, OBD-II, BeamNG.drive, Aggressive Driving Behavior}}

  \cleardoublepage

  \begin{spacing}{1.5}
    \setlength{\parskip}{0pt} % Daftar isi tetap rapat
    % Daftar isi
    \renewcommand*\contentsname{DAFTAR ISI}
    \addcontentsline{toc}{chapter}{\contentsname}
    \tableofcontents
    \cleardoublepage

    % Daftar gambar
    \renewcommand*\listfigurename{DAFTAR GAMBAR}
    \addcontentsline{toc}{chapter}{\listfigurename}
    \listoffigures
    \cleardoublepage

    % Daftar tabel
    \renewcommand*\listtablename{DAFTAR TABEL}
    \addcontentsline{toc}{chapter}{\listtablename}
    \listoftables
    \cleardoublepage
  \end{spacing}

  % Nomor halaman isi dimulai dari sini
  \pagenumbering{arabic}

  % Konten pendahuluan
  % !TEX root = ../main.tex
\chapter{PENDAHULUAN}

\section{Latar Belakang}

Jalan raya merupakan fasilitas umum yang dipakai masyarakat untuk bertransportasi setiap hari. Kelalaian pengemudi di jalan raya dapat merenggut nyawa dan menimbulkan kerugian materi. Laporan \textit{World Health Organization} (WHO) menyatakan bahwa sekitar 1,16 juta orang meninggal setiap tahun akibat kecelakaan lalu lintas di seluruh dunia \cite{WHO2026}. Di Indonesia, Pusat Informasi Kriminal Nasional (Pusiknas) Bareskrim Polri mencatat 148.798 kecelakaan lalu lintas sepanjang tahun 2023 \cite{Databoks2024}. Jumlah produksi kendaraan juga terus meningkat dalam beberapa dekade terakhir \cite{Chauhan2020}. Semakin banyak kendaraan yang berbagi ruas jalan yang sama, semakin besar pula peluang terjadinya kecelakaan.

Kesalahan manusia menjadi penyebab utama kecelakaan tersebut, yaitu 140.812 kasus atau 94,63\% dari seluruh kecelakaan pada tahun 2023 \cite{Databoks2024}. Faktor manusia yang paling sering terjadi antara lain lengah, mengantuk, tidak menjaga jarak, dan melaju dengan kecepatan tinggi \cite{Mubalus2023}. Perilaku ugal-ugalan mencakup kesalahan penilaian, pengoperasian kendaraan yang tidak tepat, sampai pengabaian rambu dan marka jalan \cite{Chauhan2020}. Perilaku ini tidak hanya membahayakan pengemudi, tetapi juga penumpang dan pengguna jalan lain \cite{Sutikno2025}. Proposal ini memakai istilah ugal-ugalan untuk perilaku mengemudi agresif yang dijelaskan pada Subbab~\ref{subsec:agresif}. Oleh karena itu, dibutuhkan sistem pemantauan yang bekerja secara \textit{real-time} untuk mendeteksi perilaku mengemudi berbahaya \cite{Heo2020}.

Kebutuhan ini muncul pada pengelola armada kendaraan komersial, seperti taksi daring, yang perlu mengetahui mitra pengemudi yang berkendara secara berisiko. Sistem manajemen armada yang membaca data kendaraan melalui port OBD-II sudah dikembangkan untuk memantau kendaraan dari jarak jauh \cite{Malekian2017}. Di Indonesia, sistem pemantau pengemudi mobil rental berbasis OBD-II menilai perilaku pengemudi dengan aturan yang ditentukan sebelumnya \cite{Rahayu2021}. Menurut tinjauan Michailidis dkk., sistem OBD-II tradisional hanya menganalisis perilaku pengemudi dari metrik dasar seperti pelanggaran kecepatan, sehingga tidak dapat menangkap pola seperti akselerasi agresif \cite{Michailidis2025}. Tinjauan yang sama mencatat bahwa model \textit{deep learning} seperti RNN, LSTM, dan CNN mampu belajar dari data sensor berurutan yang dihasilkan OBD-II \cite{Michailidis2025}.

Penelitian sebelumnya mendeteksi perilaku berbahaya dengan beberapa metode. Salah satunya menganalisis lintasan kendaraan dari rekaman kamera \cite{Sutikno2025}. Kendaraan yang berpindah jalur berkali-kali dalam waktu singkat membentuk lintasan menyerupai huruf S, dan lintasan tersebut menjadi indikator perilaku ugal-ugalan \cite{Chauhan2020}. Namun, model deteksi berbasis citra membutuhkan sumber daya komputasi yang besar ketika dijalankan pada perangkat tertanam \cite{Heo2020}. Metode ini juga hanya dapat memantau ruas jalan yang terpasang kamera, padahal pemasangan kamera pengawas di seluruh kota membutuhkan biaya besar \cite{Chauhan2020}. Penelitian lain memakai akselerometer dengan aturan dan ambang batas yang ditentukan secara manual \cite{Boonmee2009}. Pendekatan berbasis ambang juga dipakai pada perangkat IMU berdaya komputasi rendah \cite{Khan2024}. Pramunanto dkk. mengembangkan prototipe pemantau perilaku pengemudi yang mengklasifikasi kejadian longitudinal dan lateral ke dalam kelas defensif, normal, dan agresif dengan metode Na\"{\i}ve Bayes \cite{Pramunanto2019}.

Data kendaraan dapat menjadi alternatif dari pendekatan berbasis kamera. Akselerasi longitudinal dan lateral sudah cukup untuk mengenali pengereman mendadak, belokan tajam, dan akselerasi tiba-tiba dari kondisi berhenti \cite{Khan2024}. Telemetri kendaraan merekam masukan pengemudi, seperti kedalaman pedal gas, sekaligus respons kendaraan, seperti kecepatan dan putaran mesin. Data ini berbentuk deret waktu, dan \textit{Recurrent Neural Network} (RNN) sudah dipakai untuk menganalisis deret waktu kecepatan kendaraan dalam mendeteksi perilaku ugal-ugalan \cite{Thakuri2025}.

\textit{Convolutional Neural Network} (CNN) mengekstrak fitur dari data mentah tanpa perancangan fitur secara manual \cite{LeCun1998}. Kemampuan ini cocok untuk mengenali pola sesaat, seperti lonjakan akselerasi atau penekanan pedal gas yang mendadak. Perilaku ugal-ugalan juga berupa rangkaian kejadian dalam beberapa detik, misalnya akselerasi keras yang diikuti pengereman. \textit{Long Short-Term Memory} (LSTM) dirancang untuk mempelajari ketergantungan seperti ini dalam data deret waktu \cite{Hochreiter1997}. Gabungan LSTM dan CNN telah dipakai untuk mengenali perilaku mengemudi abnormal dari data kendaraan nyata \cite{Jia2020}. Cura dkk. melatih LSTM dan 1D-CNN pada data CAN bus kendaraan dan mendapatkan bahwa 1D-CNN lebih baik dalam mengenali perilaku agresif \cite{Cura2021}. Model LSTM, GRU, dan BiLSTM juga telah dipakai untuk mendeteksi perilaku mengemudi tidak aman dari data OBD-II \cite{Michailidis2025}. Oleh karena itu, penelitian ini menggabungkan CNN dengan \textit{Bidirectional} LSTM (BiLSTM) dalam satu arsitektur hibrida.

Fungsi aktivasi menentukan cara setiap neuron meneruskan sinyal. \textit{Rectified Linear Unit} (ReLU) menolkan keluaran neuron yang jumlah berbobotnya negatif, sehingga sebagian neuron dapat berhenti belajar. Kondisi ini disebut \textit{dying ReLU} \cite{Maas2013}. \textit{Parametric Rectified Linear Unit} (PReLU) mempertahankan kemiringan kecil pada sisi negatif dan mempelajari besar kemiringan tersebut selama pelatihan \cite{He2015}. Penelitian ini memakai PReLU agar neuron tetap menerima sinyal pembelajaran ketika jumlah berbobotnya negatif.

Dataset OBD-II yang terbatas dan kurang beragam menjadi salah satu tantangan penerapan \textit{machine learning} pada data kendaraan \cite{Michailidis2025}. Merekam perilaku ugal-ugalan di jalan raya juga membahayakan pengemudi dan pengguna jalan lain. Oleh karena itu, penelitian ini memakai simulator BeamNG.drive untuk mengumpulkan data latih, karena simulator tersebut mengirim data kendaraan melalui protokol UDP bawaan \cite{BeamNGProtocols}. Model yang telah dilatih kemudian diuji pada mobil nyata dengan membaca data kendaraan dari port OBD-II memakai adaptor ELM327 \cite{ELM327}. Pengujian di mobil dilakukan di area tertutup dengan skenario agresif yang terkendali. Fitur model dibatasi pada data yang tersedia di simulator dan di mobil, yaitu kecepatan, putaran mesin, posisi \textit{throttle}, dan akselerasi longitudinal yang dihitung dari selisih kecepatan.

\section{Rumusan Masalah}

Berdasarkan latar belakang yang telah diuraikan, rumusan masalah dalam Tugas Akhir ini adalah sebagai berikut:
\begin{enumerate}
    \item Metode deteksi perilaku ugal-ugalan berbasis citra membutuhkan komputasi yang besar. Metode tersebut juga hanya dapat memantau ruas jalan yang terpasang kamera. Metode berbasis data kendaraan masih menilai perilaku dengan aturan atau ambang batas yang ditentukan sebelumnya. Ambang batas tersebut menilai setiap peristiwa secara terpisah, padahal perilaku ugal-ugalan berupa rangkaian kejadian dalam beberapa detik. Oleh karena itu, diperlukan rancangan dan implementasi arsitektur hibrida CNN-BiLSTM dengan fungsi aktivasi PReLU yang mengklasifikasi perilaku pengemudi langsung dari telemetri kendaraan yang tersedia melalui OBD-II.
    \item Setelah normalisasi \textit{Z-score}, data telemetri tersebar di sekitar nol, sehingga banyak neuron pada lapisan konvolusi menerima jumlah berbobot yang negatif. ReLU menolkan keluaran neuron tersebut sehingga gradiennya hilang dan sebagian neuron dapat berhenti belajar sehingga terjadi \textit{dying ReLU}. Akurasi dan kinerja model CNN-BiLSTM dengan PReLU dalam mengklasifikasi perilaku pengemudi normal dan ugal-ugalan pada data simulator \textit{BeamNG.drive} belum diketahui.
    \item Data latih berasal dari simulator, sedangkan pengujian dilakukan pada mobil nyata melalui adaptor ELM327. Kedua sumber data berbeda dalam laju cuplik, resolusi nilai, dan jumlah besaran yang tersedia. Kemampuan model yang dilatih dengan data simulator untuk mengklasifikasi perilaku pengemudi pada data OBD-II dari mobil nyata belum diketahui.
\end{enumerate}

\section{Batasan Masalah atau Ruang Lingkup}
Penelitian ini dibatasi pada hal-hal berikut:
\begin{enumerate}
    \item Penelitian ini mengklasifikasi perilaku pengemudi ke dalam dua kelas, yaitu normal dan ugal-ugalan.
    \item Kelas ugal-ugalan didefinisikan secara operasional melalui skenario rekaman, yaitu akselerasi dan pengereman mendadak, penahanan gigi rendah, belokan tajam, dan perpindahan jalur berulang. Penelitian ini tidak mengukur niat atau kondisi emosi pengemudi.
    \item Data latih berasal dari simulator BeamNG.drive melalui protokol UDP OutGauge, dengan MotionSim sebagai data acuan untuk memeriksa akselerasi. Simulator dipilih karena alasan keamanan dan karena simulator menyediakan data fisika yang realistis pada skenario berbahaya.
    \item Pengujian pada kendaraan nyata memakai satu mobil penumpang yang dapat menjadi kendaraan komersial seperti taksi. Data dibaca melalui adaptor ELM327 dengan koneksi WiFi, dan pengujian dilakukan di area tertutup dengan skenario agresif yang terkendali.
    \item Model yang diimplementasikan adalah arsitektur hibrida CNN-BiLSTM dengan fungsi aktivasi PReLU. Masukan model berupa empat fitur inti, yaitu kecepatan, putaran mesin (RPM), posisi \textit{throttle}, dan akselerasi longitudinal yang dihitung dari selisih kecepatan. Besaran lain hanya ditambahkan jika tersedia di simulator dan pada PID mobil uji.
    \item Penelitian ini tidak memakai IMU, sehingga manuver lateral seperti zig-zag dan perpindahan jalur tidak teramati secara langsung.
    \item Data direkam oleh peneliti dan dua sampai tiga pengemudi lain dengan satu jenis transmisi. Hasil penelitian tidak mewakili seluruh populasi pengemudi.
    \item Faktor eksternal di simulator, seperti cuaca, kepadatan lalu lintas, rambu, dan waktu, dibuat konstan agar evaluasi berfokus pada kemampuan algoritma membedakan gaya mengemudi normal dan ugal-ugalan.
\end{enumerate}

\section{Tujuan}

Tujuan dari penelitian ini adalah sebagai berikut:
\begin{enumerate}
    \item Merancang dan mengimplementasikan arsitektur hibrida CNN-BiLSTM dengan aktivasi PReLU untuk mengklasifikasi perilaku pengemudi normal dan ugal-ugalan dari telemetri kendaraan.
    \item Mengukur akurasi dan kinerja model CNN-BiLSTM dengan PReLU dalam mengklasifikasi perilaku pengemudi normal dan ugal-ugalan pada data simulator \textit{BeamNG.drive}.
    \item Menguji model yang dilatih dengan data simulator pada mobil nyata melalui adaptor OBD-II ELM327.
\end{enumerate}

\section{Manfaat}

Manfaat yang diharapkan dari penelitian Tugas Akhir ini adalah sebagai berikut:

\subsection{Manfaat Praktis}
\begin{enumerate}[label=\arabic*.]
  \item \textbf{Meningkatkan Keselamatan Berkendara}: Menjadi dasar pengembangan sistem pemantauan yang memberi peringatan dini dan laporan tentang pengemudi yang berkendara secara berisiko.
  \item \textbf{Penilaian Risiko Pengemudi}: Memberi alternatif bagi perusahaan asuransi kendaraan dan pengelola armada komersial untuk menilai profil risiko pengemudi dari data OBD-II. Pemantauan armada berbasis data OBD-II sudah diterapkan untuk menghitung jarak tempuh dan konsumsi bahan bakar \cite{Malekian2017}.
  \item \textbf{Riset yang Aman dan Hemat Biaya}: Menunjukkan pemanfaatan simulator \textit{high-fidelity} seperti BeamNG.drive untuk mengumpulkan data latih, sehingga pengujian dengan mobil nyata hanya diperlukan pada tahap akhir dan di area yang aman.
\end{enumerate}

\subsection{Manfaat Teoretis}
\begin{enumerate}[label=\arabic*.]
  \item \textbf{Pengembangan Arsitektur \textit{Deep Learning}}: Menambah kajian tentang penerapan fungsi aktivasi PReLU pada arsitektur hibrida CNN-BiLSTM untuk data telemetri kendaraan.
  \item \textbf{Dari Simulator ke Kendaraan Nyata}: Memberi gambaran tentang kemampuan model yang dilatih dengan data simulator ketika diterapkan pada data OBD-II dari mobil nyata.
  \item \textbf{Referensi Akademis}: Menjadi rujukan bagi penelitian selanjutnya di bidang teknik komputer, khususnya sistem transportasi cerdas dan pengolahan data sensor dengan \textit{Deep Learning}.
\end{enumerate}

  \cleardoublepage

  % Konten tinjauan pustaka
  % !TEX root = ../main.tex
\chapter{TINJAUAN PUSTAKA}

\section{Hasil penelitian/perancangan terdahulu}
Penelitian tentang deteksi perilaku mengemudi berbahaya memakai dua pendekatan utama, yaitu pendekatan berbasis citra atau video dan pendekatan berbasis data kendaraan. Bagian ini menguraikan lima belas penelitian terdahulu yang menjadi acuan penelitian ini, diurutkan berdasarkan tahun publikasi.

Boonmee dan Tangamchit mengembangkan sistem portabel untuk mendeteksi perilaku mengemudi buruk pada bus menggunakan akselerometer dua sumbu. Sistem tersebut mengenali empat peristiwa berkendara, yaitu \textit{take off}, \textit{brake}, \textit{turn left}, dan \textit{turn right}, dengan metode berbasis aturan (\textit{rule-based}). Setiap peristiwa kemudian diberi skor dengan aturan \textit{fuzzy logic}. Hasil pengujian yang dibandingkan dengan penilaian 11 relawan menunjukkan akurasi 88,04\% untuk \textit{take off}, 87,21\% untuk \textit{brake}, 88,84\% untuk \textit{turn right}, dan 84,36\% untuk \textit{turn left}.

Carmona dkk. mengembangkan alat analisis perilaku pengemudi berbiaya rendah yang menggabungkan data sensor kendaraan dari CAN bus, \textit{inertial measurement unit} (IMU), dan GPS. Sistem tersebut mengolah sejumlah parameter dan sinyal pada domain waktu dan frekuensi untuk mendeteksi perilaku agresif secara \textit{real-time}. Lokalisasi GPS menambahkan konteks berupa informasi peta digital dan batas kecepatan. Sistem tersebut diuji pada skenario perkotaan, antarkota, dan jalan tol.

Pramunanto dkk. mengembangkan prototipe sistem pemantau perilaku pengemudi kendaraan roda empat dengan metode Na\"{\i}ve Bayes. Sistem tersebut memakai deteksi kejadian dan ekstraksi fitur untuk mengklasifikasi perilaku pengemudi ke dalam tiga kelas, yaitu defensif, normal, dan agresif. Validasi menunjukkan akurasi 86\% untuk kejadian longitudinal dan 95,8\% untuk kejadian lateral. Penelitian ini melanjutkan arah tersebut dengan mengklasifikasi perilaku per jendela waktu memakai CNN-BiLSTM.

Chauhan dkk. mengusulkan sistem \textit{deep learning} untuk mendeteksi perilaku mengemudi ugal-ugalan di jalan raya dari rekaman video. Sistem tersebut terdiri atas dua modul, yaitu ekstraksi lintasan kendaraan dan deteksi penyok (\textit{dent}) pada badan kendaraan. Lintasan kendaraan diperoleh dari koordinat objek hasil deteksi YOLOv3. Sudut antarvektor pergerakan pada tiga \textit{frame} berurutan dipakai untuk memperkirakan probabilitas perilaku ugal-ugalan. Hasil kedua modul digabungkan, dan sistem divalidasi dengan video dari YouTube.

Heo dkk. merancang sistem tertanam untuk mendeteksi perilaku mengemudi ugal-ugalan secara otonom dan \textit{real-time} menggunakan \textit{embedded GPU}, kamera, dan model \textit{deep learning} ringan. Sistem tersebut membentuk lintasan gerak kendaraan dari citra per \textit{frame} dan memanfaatkan kisi virtual (\textit{virtual grid}) pada jalan agar pelacakan tetap akurat pada video berlaju \textit{frame} rendah, yaitu 5 \textit{fps}. Evaluasi pada lintasan kendaraan nyata dan lintasan mobil kendali jarak jauh menghasilkan \textit{recall} 100\% dan presisi 67\%. Keterbatasan sumber daya komputasi pada platform tertanam menjadi tantangan utama yang dibahas pada penelitian tersebut.

Jia dkk. mengusulkan model pengenalan perilaku mengemudi abnormal yang menggabungkan \textit{Long Short-Term Memory} dan \textit{Convolutional Neural Network} (LSTM-CNN). Perilaku abnormal yang dimaksud antara lain akselerasi cepat, pengereman darurat, dan perpindahan jalur yang cepat. Menurut penelitian tersebut, sebagian besar penelitian sebelumnya memakai metode \textit{clustering} dan \textit{shallow learning} yang sulit mengenali jenis perilaku abnormal secara akurat. Jia dkk. membangun dataset dari data mengemudi kendaraan nyata dengan mendeteksi titik akselerasi dan deselerasi ekstrem melalui analisis statistik, dan model LSTM-CNN yang dilatih dengan dataset tersebut memberikan hasil yang lebih baik. Berbeda dengan penelitian tersebut, penelitian ini memakai data telemetri dari simulator dan menempatkan lapisan konvolusi sebelum lapisan BiLSTM.

Cura dkk. mengembangkan model LSTM dan CNN untuk mengklasifikasi perilaku pengemudi bus berdasarkan deselerasi, putaran mesin saat menginjak pedal, belokan, dan perpindahan jalur. Uji deselerasi, putaran mesin, dan belokan dilakukan di lintasan uji beton, sedangkan uji pindah jalur dilakukan di jalan tol beraspal. Model hanya memakai aliran data dari CAN bus kendaraan yang diolah menjadi parameter karakteristik. LSTM dan 1D-CNN menghasilkan tingkat keberhasilan yang setara, dan CNN lebih baik dalam mengenali perilaku agresif.

Rahayu membangun sistem pemantau untuk layanan rental mobil yang membaca RPM, kecepatan, beban mesin, dan suhu mesin dari OBD-II melalui Raspberry Pi 2, serta lokasi GPS dari smartphone. Sistem tersebut menilai perilaku pengemudi dengan aturan yang ditentukan dan menghasilkan tingkat kesalahan mengemudi. Pengujian menghasilkan 173 data dengan 9 pelanggaran aturan, atau tingkat kesalahan mengemudi 5,20\%.

Khan dkk. membangun perangkat \textit{plug-and-play} untuk menganalisis perilaku mengemudi secara \textit{real-time} menggunakan prosesor berdaya komputasi rendah dan sensor \textit{inertial measurement unit} (IMU) berbiaya rendah. Penelitian tersebut tidak memakai jaringan saraf tiruan karena kompleksitas komputasi dan mahalnya prosesor yang dibutuhkan. Algoritma yang diusulkan memantau percepatan longitudinal dan lateral untuk mendeteksi pengereman mendadak, belokan tajam, dan akselerasi tiba-tiba dari kondisi berhenti. Pengujian pada \textit{test bench} menunjukkan bahwa pendekatan algoritmik tersebut memiliki latensi yang cukup rendah untuk analisis perilaku mengemudi secara \textit{real-time}.

Muhamad Fadzil dkk. mengklasifikasi perilaku pengemudi dari dataset OBD-II daring yang diperoleh dari Kaggle. Fitur yang dipilih adalah putaran mesin, kecepatan kendaraan, posisi \textit{throttle}, dan beban mesin. Label aman atau agresif dibuat dengan algoritma K-Means, lalu beberapa model \textit{machine learning} dilatih untuk memprediksi label tersebut. Decision Tree mencapai akurasi rata-rata 98,72\%, dan AdaBoost menaikkannya menjadi 99,48\%. Label pada penelitian tersebut berasal dari pengelompokan data, sedangkan penelitian ini memberi label berdasarkan skenario rekaman.

Singh dan Sharma memakai data OBD-II dan pengukuran gaya-g untuk mengklasifikasi sepuluh perilaku mengemudi, antara lain \textit{weaving}, \textit{swerving}, dan \textit{sideslipping}. Model yang dipakai adalah Random Forest, AdaBoost, dan K-Nearest Neighbors, dan Random Forest mencapai akurasi tertinggi.

Thakuri dkk. mengembangkan aplikasi penghitung status ugal-ugalan (\textit{Reckless Status}) dari rekaman video jalan raya. Kendaraan dideteksi dengan YOLO dan dilacak dengan \textit{Kalman filter} untuk memperoleh data kecepatan dan lintasan. Parameter statistik kecepatan, yaitu rata-rata, standar deviasi, varians, \textit{kurtosis}, \textit{skewness}, dan \textit{peak-to-peak}, kemudian dianalisis sebagai data deret waktu dengan \textit{Recurrent Neural Network} (RNN). Model RNN tersebut mencapai akurasi 94,12\%, presisi 1,0, \textit{recall} 91,6\%, dan \textit{F1-score} 0,8979. Ambang batas yang dipakai ditentukan secara heuristik dan hanya sesuai untuk video percobaan yang direkam.

Sutikno dkk. mengembangkan sistem deteksi pengemudi ugal-ugalan berdasarkan lintasan kendaraan yang terdiri atas tiga subsistem, yaitu deteksi mobil, pelacakan mobil dengan ByteTrack, dan klasifikasi lintasan. Untuk klasifikasi lintasan, penelitian tersebut mengusulkan YOLO11n-cls yang ditingkatkan dengan blok konvolusi dan C3k2 tambahan. Model tersebut mengklasifikasikan lintasan ke dalam tiga kelas, yaitu berkendara lurus, pindah jalur halus, dan ugal-ugalan. Model usulan mencapai akurasi 0,935 dengan waktu inferensi 0,5 ms, meningkat 4,4\% dibandingkan YOLO11n-cls. Model tersebut masih lemah dalam memprediksi kelas pindah jalur halus.

Michailidis dkk. meninjau aplikasi \textit{machine learning} berbasis data OBD-II untuk efisiensi energi, pengendalian emisi, analisis perilaku pengemudi, deteksi anomali, dan keamanan siber. Menurut tinjauan tersebut, aplikasi OBD-II konvensional memakai aturan atau ambang tetap dan hanya mengenali masalah yang sudah ditentukan, sedangkan model RNN, LSTM, dan CNN mampu belajar dari data sensor berurutan. Tinjauan tersebut merangkum beberapa penelitian perilaku pengemudi berbasis OBD-II. Lee dkk. melatih LSTM, GRU, BiLSTM, dan BiGRU pada data OBD-II dari enam pengemudi di lintasan terkendali untuk empat perilaku, yaitu aman, akselerasi agresif, belokan tajam, dan pengereman keras, dan GRU mencapai akurasi 97,5\%. Song dkk. mengenali gaya mengemudi agresif, normal, dan konservatif dengan GMM semi-supervised, dan jendela 30 detik memberi hasil terbaik dengan akurasi 92\%. Tinjauan tersebut juga mencatat tantangan penerapan, yaitu dataset yang terbatas dan kurang beragam, data OBD-II yang berderau dengan laju data rendah, model \textit{deep learning} yang sering terlalu berat untuk perangkat kendaraan, serta PID khusus merek yang menyulitkan pemindahan model. Penelitian ini menjawab tantangan tersebut dengan data latih dari simulator, model berukuran kecil, dan fitur yang hanya memakai PID standar.

Sa\u{g}ba\c{s} mengembangkan kerangka kerja berbiaya rendah yang merekam seluruh data melalui aplikasi smartphone yang terhubung ke adaptor OBD-II Bluetooth. Data yang digabungkan meliputi data unit kendali mesin (ECU), sensor gerak ponsel, metadata jalan, dan sinyal audio. Perilaku mengemudi dikelompokkan ke dalam tiga kelas, yaitu tenang, normal, dan agresif, dengan jendela waktu 3, 5, 7, dan 9 detik. CNN mencapai kinerja tertinggi, dan penggabungan beberapa sumber data menaikkan akurasi klasifikasi.

Dari penelitian-penelitian tersebut, penelitian berbasis data kendaraan sebagian besar masih menilai perilaku dengan aturan atau ambang tetap, atau memberi label melalui pengelompokan data. Penelitian ini melatih model CNN-BiLSTM dengan PReLU memakai data simulator yang labelnya ditentukan oleh skenario rekaman, lalu menguji model tersebut pada data OBD-II dari mobil nyata.

\section{Teori/Konsep Dasar}
\subsection{Perilaku Mengemudi dan Keselamatan Lalu Lintas}
\subsubsection{Kecelakaan Lalu Lintas dan Faktor Penyebabnya}
Menurut WHO, kecelakaan lalu lintas masih menjadi masalah keselamatan global yang merenggut 1,16 juta jiwa setiap tahun. Mubalus membagi faktor penyebab kecelakaan lalu lintas ke dalam tiga kelompok, yaitu faktor manusia, faktor jalan, dan faktor kendaraan. Dari ketiganya, kesalahan manusia menjadi penyebab yang paling sering terjadi. Mubalus mengutip bahwa hampir 86\% kecelakaan lalu lintas disebabkan oleh kelalaian manusia, dengan faktor tersering berupa lengah, mengantuk, tidak menjaga jarak, dan melaju dengan kecepatan tinggi. Dari sisi hukum, mengemudikan kendaraan secara tidak wajar dan tanpa konsentrasi diatur dalam Pasal~283 UU No.~22 Tahun 2009. Kecepatan juga merupakan faktor risiko yang terukur. WHO menyatakan bahwa setiap kenaikan 1\% kecepatan rata-rata berkaitan dengan kenaikan 4\% risiko kecelakaan fatal.

\subsubsection{Perilaku Mengemudi Agresif/Ugal-ugalan}
\label{subsec:agresif}
Handayani dkk., yang mengutip Tasca (2000), mendefinisikan mengemudi agresif (\textit{aggressive driving}) sebagai perilaku sengaja yang cenderung meningkatkan risiko kecelakaan dan didorong oleh ketidaksabaran, kekesalan, permusuhan, atau upaya menghemat waktu. Contoh perilaku tersebut adalah berkendara secara zig-zag dan membuntuti kendaraan lain dengan jarak terlalu dekat. Thakuri dkk. mengakui bahwa istilah \textit{reckless} bersifat abstrak dan sulit didefinisikan, sehingga mereka mendefinisikannya melalui kecepatan dan parameter statistik dengan ambang batas heuristik. Penelitian ini memakai istilah ugal-ugalan untuk perilaku mengemudi agresif menurut definisi Tasca. Kedua istilah tersebut merujuk ke kelas yang sama. Definisi Tasca memuat unsur niat dan emosi yang tidak dapat diukur dari data kendaraan. Oleh karena itu, penelitian ini memakai definisi operasional: kelas ugal-ugalan adalah sesi mengemudi dengan akselerasi dan pengereman mendadak, penahanan gigi rendah, belokan tajam, dan perpindahan jalur berulang. Model mengenali perilaku agresif yang teramati dari telemetri, bukan niat pengemudi.

\subsubsection{Indikator Perilaku Ugal-ugalan}
Boonmee dan Tangamchit memakai akselerasi sumbu X untuk mengenali \textit{take off} dan \textit{brake}, serta akselerasi sumbu Y untuk mengenali gerakan ke kiri dan ke kanan. Mereka menetapkan batas kondisi tanpa aksi, yaitu akselerasi di bawah 87,5~mg pada sumbu X dan di bawah 150~mg pada sumbu Y. Satuan mg (\textit{milli-g}) berarti seperseribu percepatan gravitasi, sehingga 150~mg sama dengan sekitar 1,5~m/s$^{2}$. Mereka juga memakai \textit{connected time} untuk membedakan peristiwa kuat dari peristiwa ringan. Khan dkk. memakai akselerasi lateral untuk mengenali belokan tajam dan perpindahan jalur mendadak.

Chauhan dkk. menggolongkan kendaraan yang berpindah jalur berkali-kali dalam waktu singkat hingga membentuk lintasan menyerupai huruf S sebagai cara berkendara yang berbahaya. Perubahan kecepatan yang mendadak juga menjadi indikator perilaku ugal-ugalan, yaitu akselerasi dan deselerasi kuat yang terjadi dalam waktu berdekatan. Thakuri dkk. memakai parameter \textit{peak-to-peak} kecepatan, yaitu selisih nilai tertinggi dan terendah, untuk membedakan cara berkendara yang stabil dari cara mengemudi yang agresif. Indikator lateral seperti belokan tajam dan perpindahan jalur memerlukan sensor gerak. Singh dan Sharma menambahkan pengukuran gaya-g, dan Carmona dkk. menambahkan IMU di samping data kendaraan.

\subsection{Telemetri Kendaraan}
\subsubsection{Pengertian dan Parameter Telemetri}
Telemetri kendaraan adalah proses pengukuran dan pengiriman data parameter kendaraan secara otomatis ke penerima untuk dipantau dan dianalisis. Parameter yang umum diukur adalah kecepatan, akselerasi longitudinal dan lateral, laju rotasi kendaraan, putaran mesin, posisi gigi transmisi, serta posisi pedal gas dan pedal rem. Pada simulator, program membaca parameter-parameter tersebut langsung dari mesin fisika, sehingga tidak dibutuhkan sensor fisik. Pada mobil nyata, sebagian parameter tersebut dapat dibaca melalui port OBD-II yang dijelaskan pada Subbab~\ref{subsec:obd}.

Besaran gerak kendaraan dinyatakan pada sistem koordinat yang menempel pada badan kendaraan. Menurut Rajamani, sumbu longitudinal menunjuk ke arah depan kendaraan, sumbu lateral menunjuk ke samping, dan sumbu vertikal menunjuk ke atas. Dengan sistem koordinat ini, akselerasi longitudinal menggambarkan percepatan dan pengereman, sedangkan akselerasi lateral menggambarkan gaya samping ketika kendaraan berbelok. Kedua nilai tersebut tidak bergantung pada arah hadap kendaraan di peta.

\subsubsection{Kecepatan, Akselerasi, dan Deselerasi}
Kecepatan adalah perubahan posisi kendaraan per satuan waktu, sedangkan akselerasi adalah perubahan kecepatan per satuan waktu. Untuk dua cuplikan berurutan dengan selang waktu $\Delta t$, keduanya dihitung dengan
\begin{equation}
  v = \frac{\Delta x}{\Delta t}
  \label{eq:kecepatan}
\end{equation}
\begin{equation}
  a = \frac{\Delta v}{\Delta t}
  \label{eq:akselerasi}
\end{equation}
dengan $\Delta x$ adalah perubahan posisi dalam meter, $\Delta v$ adalah perubahan kecepatan dalam meter per detik (m/s), dan $\Delta t$ adalah selang waktu dalam detik. Satuan akselerasi adalah meter per detik kuadrat (m/s$^{2}$). Akselerasi positif berarti kendaraan bertambah cepat. Akselerasi negatif disebut deselerasi, yaitu kondisi saat kendaraan melambat, misalnya ketika pengemudi mengerem.

Pada penelitian ini, akselerasi longitudinal dihitung dari selisih kecepatan, baik pada simulator maupun pada mobil, karena adaptor OBD-II yang dipakai tidak menyediakan data akselerasi. Cara ini memperbesar derau, karena kesalahan kecil pada kecepatan dibagi dengan $\Delta t$ yang kecil. Menurut standar SAE J1979, kecepatan kendaraan pada OBD-II dikirim dalam satuan km/jam bulat. Perubahan satu satuan, yaitu 1~km/jam atau sekitar 0,28~m/s, pada $\Delta t = 0{,}2$~s sudah menghasilkan lonjakan akselerasi sekitar 1,4~m/s$^{2}$. Oleh karena itu, akselerasi hasil selisih diratakan dengan rata-rata bergerak (\textit{moving average}) atas $m$ cuplikan terakhir:
\begin{equation}
  \bar{a}_{k} = \frac{a_{k} + a_{k-1} + \dots + a_{k-m+1}}{m}
  \label{eq:moving-average}
\end{equation}
dengan $a_{k}$ adalah akselerasi hasil selisih pada cuplikan ke-$k$. Sebagai contoh, dengan $m = 3$, nilai $\bar{a}_{k}$ adalah rata-rata akselerasi dari cuplikan saat ini dan dua cuplikan sebelumnya. Kecepatan dari simulator juga dibulatkan ke km/jam bulat sebelum akselerasi dihitung, sehingga derau pada data latih menyerupai derau pada data mobil.

\subsubsection{Putaran Mesin (RPM) dan Posisi Gigi}
Putaran mesin dalam satuan RPM adalah jumlah putaran poros engkol dalam satu menit. Nilai RPM naik ketika pengemudi menekan pedal gas lebih dalam. Pada kecepatan kendaraan yang sama, RPM ditentukan oleh posisi gigi transmisi. Gigi rendah menghasilkan RPM tinggi dan tenaga dorong besar, sedangkan gigi tinggi menghasilkan RPM rendah.

Pengemudi agresif cenderung menekan pedal gas dalam sehingga mesin berputar tinggi dan transmisi tertahan pada gigi rendah. Gaya mengemudi yang halus cenderung berada pada gigi tinggi dengan RPM rendah. Hasil pendahuluan pada BAB IV menunjukkan bahwa RPM dan posisi gigi menjadi pembeda kelas terkuat. Muhamad Fadzil dkk. juga memilih putaran mesin, kecepatan, dan posisi \textit{throttle} sebagai fitur klasifikasi perilaku pengemudi dari data OBD-II. Posisi gigi hanya dapat dibaca jika mobil mendukung PID yang sesuai, sehingga penelitian ini tidak memakai posisi gigi sebagai fitur inti.

\subsubsection{Kedalaman Pedal Gas dan Pedal Rem}
Kedalaman pedal gas (\textit{throttle}) menyatakan seberapa dalam pengemudi menekan pedal akselerator, dan kedalaman pedal rem (\textit{brake}) menyatakan seberapa dalam pengemudi menekan pedal rem. Penelitian ini menyimpan kedua nilai dalam persen, dengan 0\% berarti pedal tidak ditekan dan 100\% berarti pedal ditekan penuh.

Kecepatan dan akselerasi menunjukkan hasil gerak kendaraan, sedangkan kedua pedal menunjukkan masukan dari pengemudi. Pengemudi ugal-ugalan cenderung menekan pedal gas dalam, melepasnya dengan cepat, lalu menekan pedal rem secara mendadak. Penelitian ini memakai kedalaman pedal gas sebagai fitur inti, karena besaran yang setara dapat dibaca melalui OBD-II. Kedalaman pedal rem hanya dipakai jika mobil uji menyediakannya.


\subsubsection{Telemetri pada Kendaraan Nyata melalui OBD-II}
\label{subsec:obd}
\textit{On-Board Diagnostics} II (OBD-II) adalah antarmuka diagnostik pada kendaraan yang memberi akses ke data unit kendali mesin (ECU). Permintaan data OBD-II terdiri atas byte mode dan nomor parameter (\textit{parameter identification}, PID). Mode dan PID tersebut dijelaskan dalam standar SAE J1979, dan produsen kendaraan juga dapat mendefinisikan PID sendiri. Kendaraan tidak wajib mendukung semua PID. Menurut lembar data ELM327, PID 00 pada mode 01 menunjukkan daftar PID yang didukung kendaraan dan wajib didukung oleh semua kendaraan. PID pada mode 22 umumnya milik produsen dan tidak dipublikasikan secara luas, sehingga isinya berbeda untuk setiap merek kendaraan.

ELM327 adalah penerjemah yang menjembatani port OBD-II dengan antarmuka serial standar, sehingga komputer dapat membaca data kendaraan. ELM327 dapat mendeteksi dan menerjemahkan sembilan protokol OBD secara otomatis. Penelitian ini memakai adaptor ELM327 dengan koneksi WiFi, sehingga program Python di laptop membaca data melalui jaringan nirkabel. Setiap PID diminta satu per satu, sehingga laju cuplik bergantung pada jumlah PID yang dibaca dan waktu respons kendaraan. Menurut Michailidis dkk., laju cuplik OBD-II umumnya berkisar 1 sampai 100~Hz, bergantung pada jumlah PID yang diminta, batasan protokol, dan beban ECU. Lembar data ELM327 juga mengingatkan bahwa kendaraan lama mungkin tidak dapat menerima permintaan yang terlalu cepat.

Penelitian terdahulu memakai beberapa cara untuk membaca data kendaraan nyata. Ceven dkk. membangun perangkat akuisisi dari Raspberry Pi 3 Model B+ dan modul PiCAN 2 yang membaca data ECU secara \textit{real-time} melalui konektor OBD-II dan CAN bus. Carmona dkk. menambahkan IMU dan GPS di samping data kendaraan. Sa\u{g}ba\c{s} memakai smartphone yang terhubung ke adaptor OBD-II dan sensor gerak ponsel. Malekian dkk. membaca kecepatan dan aliran massa udara dari OBD-II untuk menghitung jarak tempuh dan konsumsi bahan bakar armada kendaraan, lalu mengirimkan datanya melalui WiFi ke server.

Penelitian ini tidak memakai IMU. Fitur model dibatasi pada besaran yang tersedia di simulator dan di mobil, sesuai Tabel~\ref{tab:padanan}.

\begin{table}[H]
  \centering
  \caption{Padanan fitur inti pada simulator dan mobil}
  \label{tab:padanan}
  \small
  \begin{tabular}{|l|l|l|p{4.2cm}|}
    \hline
    \textbf{Fitur} & \textbf{Simulator} & \textbf{Mobil (mode 01)} & \textbf{Keterangan} \\
    \hline
    Kecepatan & OutGauge & PID 0D & km/jam bulat \\
    \hline
    Putaran mesin & OutGauge & PID 0C & rpm \\
    \hline
    Posisi \textit{throttle} & OutGauge & PID 11 & PID 49 (pedal akselerator) dipakai jika tersedia \\
    \hline
    Akselerasi longitudinal & selisih kecepatan & selisih kecepatan & diratakan dengan Persamaan~\ref{eq:moving-average} \\
    \hline
  \end{tabular}
\end{table}

Posisi gigi, kedalaman pedal rem, dan sudut kemudi tidak termasuk fitur inti karena ketersediaannya berbeda pada setiap mobil. Jika pemeriksaan PID menunjukkan bahwa mobil uji menyediakan salah satu besaran tersebut, dan besaran yang sama juga tersedia di simulator, besaran itu dapat ditambahkan sebagai fitur.

\subsection{Simulator BeamNG}
\subsubsection{Gambaran Umum dan Fisika Soft-body}
BeamNG adalah simulator kendaraan yang memodelkan fisika kendaraan dengan pendekatan \textit{soft-body}. Pendekatan ini memperlakukan kendaraan sebagai benda yang dapat berubah bentuk, sedangkan model \textit{rigid body} menganggap bentuk benda tetap. Pada BeamNG, kendaraan tersusun atas titik massa (\textit{node}) yang dihubungkan oleh elemen \textit{beam}. Setiap \textit{beam} bekerja seperti pegas dan peredam. Simulator menjumlahkan gaya dari semua \textit{beam} yang terhubung ke suatu \textit{node}, lalu menghitung gerak \textit{node} tersebut dengan hukum kedua Newton.

Penelitian ini memakai BeamNG karena simulator menghitung respons kendaraan terhadap pedal dan kemudi dari model fisika. Akibatnya, akselerasi, laju yaw, dan putaran mesin yang direkam mengikuti dinamika kendaraan. Simulator juga memungkinkan peneliti merekam perilaku ugal-ugalan berulang kali tanpa risiko bagi pengemudi dan pengguna jalan.

\subsubsection{BeamNG.drive, BeamNG.tech, dan Pengembangan ADAS}
\textit{Advanced Driver Assistance Systems} (ADAS) adalah sistem yang memantau kondisi kendaraan dan lingkungan untuk membantu pengemudi, misalnya peringatan tabrakan dan pengereman darurat otomatis. Pengujian ADAS di jalan nyata mahal dan berbahaya, sehingga pengembang membutuhkan lingkungan simulasi. Sistem klasifikasi perilaku mengemudi dalam penelitian ini dapat menjadi salah satu komponen ADAS, yaitu modul yang mengenali pola mengemudi berisiko dari data kendaraan.

BeamNG tersedia dalam dua versi. BeamNG.tech adalah versi untuk riset dan industri yang menyediakan antarmuka pemrograman BeamNGpy untuk mengendalikan skenario dan membaca sensor kendaraan dari Python. BeamNG.drive adalah versi umum yang memakai mesin fisika yang sama. Penelitian ini memakai BeamNG.drive karena seluruh data yang dibutuhkan dapat diperoleh melalui protokol UDP bawaan pada subbab berikutnya, tanpa lisensi riset.

\subsubsection{Akuisisi Data Telemetri Melalui Protokol UDP}
\textit{User Datagram Protocol} (UDP) adalah protokol lapisan transport yang mengirim data dalam bentuk paket tanpa membangun koneksi terlebih dahulu. UDP tidak mengirim ulang paket yang hilang dan tidak menjamin urutan paket. Sebagai gantinya, beban tambahan UDP kecil dan latensinya rendah, sehingga protokol ini cocok untuk aliran data terus-menerus seperti telemetri.

BeamNG.drive menyediakan dua protokol UDP bawaan yang diaktifkan melalui menu \textit{Options} \textgreater{} \textit{Other} \textgreater{} \textit{Protocols}:
\begin{enumerate}
  \item \textbf{OutGauge}, yaitu protokol dengan format yang sama dengan simulator Live for Speed. Satu paket berukuran 96 \textit{byte} dan memuat informasi dasbor, antara lain kecepatan (m/s), RPM, posisi gigi, kedalaman pedal gas, pedal rem, dan kopling.
  \item \textbf{MotionSim}, yaitu protokol untuk menggerakkan simulator gerak (\textit{motion platform}). Satu paket berukuran 88 \textit{byte}, diawali penanda empat karakter \texttt{"BNG1"}, lalu diikuti 21 bilangan \textit{float}: posisi, kecepatan, akselerasi tanpa gravitasi, vektor arah atas, sudut \textit{roll}--\textit{pitch}--\textit{yaw}, kecepatan sudut, dan percepatan sudut kendaraan.
\end{enumerate}
Penelitian ini memakai OutGauge sebagai sumber fitur model, karena isinya setara dengan data yang dapat dibaca melalui OBD-II. MotionSim dipakai untuk memeriksa akselerasi hasil selisih kecepatan terhadap akselerasi yang dihitung mesin fisika simulator.

Program penerima membaca kedua aliran paket dan menggabungkannya menjadi satu baris data pada setiap detak jam dengan frekuensi tetap $f_{s}$. Selang antarcuplikan dan jumlah cuplikan dari rekaman berdurasi $T$ detik adalah
\begin{equation}
  \Delta t = \frac{1}{f_{s}}, \qquad N = f_{s}\,T
  \label{eq:fs}
\end{equation}
Pada $f_{s} = 10$~Hz, selang antarcuplikan adalah 0,1~s, dan rekaman satu menit menghasilkan 600 cuplikan. Jam tetap membuat jarak antarbaris seragam walaupun waktu tiba paket UDP berubah-ubah. Karena UDP dapat kehilangan paket, jumlah cuplikan yang tersimpan dapat lebih kecil dari $N$. Program melewati baris yang sumber datanya terlambat, dan tahap pembersihan data menangani celah yang terbentuk.

\subsection{Prapemrosesan Data Deret Waktu}
\subsubsection{Pembersihan Data dan Segmentasi}
Pembersihan data adalah proses menangani data yang tidak valid agar tidak menurunkan kualitas model. Data telemetri dari UDP dapat mengandung cuplikan yang hilang, nilai yang tidak terbaca, dan nilai ekstrem akibat kejadian di luar perilaku mengemudi, misalnya tabrakan. Baris dengan nilai kosong atau tidak hingga dihapus, dan nilai ekstrem dipotong (\textit{clipping}) pada batas fisik yang masuk akal.

Penyaringan data perlu memperhatikan hubungan antara filter dan label. Jika filter yang diterapkan pada satu kelas memakai variabel yang juga menjadi masukan model, misalnya membuang baris berkecepatan tinggi hanya dari kelas normal, model dapat membedakan kelas cukup dengan mengenali batas filter tersebut. Akurasi yang diperoleh kemudian tidak mencerminkan kemampuan mengenali pola perilaku. Penyaringan yang sah berlaku sama untuk semua kelas, misalnya membuang cuplikan ketika kendaraan diam.

Cuplikan yang hilang dapat diisi dengan interpolasi linear, yaitu menarik garis lurus di antara dua cuplikan terdekat. Interpolasi hanya wajar untuk celah yang pendek. Pada celah panjang atau baris yang dibuang oleh penyaringan, nilai hasil interpolasi tidak pernah terjadi. Pendekatan yang lebih aman adalah segmentasi, yaitu memotong rekaman menjadi segmen-segmen yang barisnya berurutan, lalu mengambil setiap jendela data hanya dari dalam satu segmen. Segmen baru dimulai ketika selisih waktu antarbaris melebihi batas tertentu atau ketika ada baris yang dibuang.

\subsubsection{Penyamaan Laju Cuplik (\textit{Resampling})}
Data simulator direkam pada 10~Hz, sedangkan laju cuplik ELM327 bergantung pada jumlah PID dan waktu respons kendaraan. Model hanya mengenali pola dengan benar jika selang antarcuplikan saat pengujian sama dengan saat pelatihan. Oleh karena itu, data simulator di-\textit{resample} ke laju cuplik mobil $f_{s}$ dengan mengambil nilai pada titik waktu berjarak $1/f_{s}$ di dalam setiap segmen. Contohnya, jika $f_{s} = 5$~Hz, satu dari setiap dua cuplikan simulator dipakai. Ukuran jendela ditetapkan dalam detik, sehingga jumlah cuplikan dalam jendela tiga detik adalah $W = 3 f_{s}$.

\subsubsection{Normalisasi (Min-Max dan Z-score)}
Fitur telemetri memiliki skala yang berbeda. RPM dapat bernilai ribuan, sedangkan laju yaw hanya bernilai sepersekian radian per detik. Jika data dimasukkan tanpa normalisasi, fitur berskala besar mendominasi perhitungan gradien, sehingga pelatihan menjadi lambat dan tidak stabil.

Normalisasi \textit{min-max} mengubah setiap nilai ke rentang 0 sampai 1:
\begin{equation}
  x' = \frac{x - x_{\min}}{x_{\max} - x_{\min}}
  \label{eq:minmax}
\end{equation}
Normalisasi \textit{Z-score} mengubah nilai agar fitur memiliki rata-rata 0 dan simpangan baku 1:
\begin{equation}
  x' = \frac{x - \mu}{\sigma}
  \label{eq:zscore}
\end{equation}
dengan $x_{\min}$ dan $x_{\max}$ adalah nilai terkecil dan terbesar fitur, $\mu$ adalah rata-rata fitur, dan $\sigma$ adalah simpangan baku fitur. Sebagai contoh, RPM 4.000 dengan rata-rata 3.000 dan simpangan baku 1.000 menjadi 1, sedangkan RPM 2.000 menjadi $-1$. Nilai di bawah rata-rata selalu menjadi negatif setelah normalisasi \textit{Z-score}.

Normalisasi \textit{min-max} sensitif terhadap satu nilai ekstrem, karena nilai tersebut menentukan $x_{\max}$ atau $x_{\min}$. Oleh karena itu, penelitian ini memakai \textit{Z-score}. Nilai $\mu$ dan $\sigma$ dihitung dari data latih saja, lalu dipakai untuk data validasi, data uji, dan inferensi \textit{real-time}, sehingga informasi data uji tidak masuk ke proses pelatihan.

\subsubsection{Sliding Window}
Satu cuplikan telemetri hanya merekam kondisi kendaraan pada satu titik waktu, sehingga belum cukup untuk menilai gaya mengemudi. Pengereman mendadak atau akselerasi agresif baru terlihat dari perubahan nilai selama beberapa detik. \textit{Sliding window} memotong rekaman telemetri menjadi potongan pendek berukuran tetap yang disebut jendela. Setiap jendela memuat $W$ cuplikan berurutan dari seluruh fitur. Setelah satu jendela terbentuk, jendela bergeser sejauh $S$ cuplikan untuk membentuk jendela berikutnya. Dari satu segmen berisi $N$ cuplikan diperoleh jendela sebanyak
\begin{equation}
  N_{w} = \left\lfloor \frac{N - W}{S} \right\rfloor + 1
  \label{eq:jumlah-window}
\end{equation}
Tanda $\lfloor \cdot \rfloor$ berarti hasil pembagian dibulatkan ke bawah. Sebagai contoh, segmen berisi $N = 200$ cuplikan dengan $W = 30$ dan $S = 5$ menghasilkan $\lfloor 170/5 \rfloor + 1 = 35$ jendela. Jika $S < W$, dua jendela yang berurutan saling tumpang tindih dan berbagi $W - S$ cuplikan.

Satu jendela menjadi satu sampel masukan bagi model dan memperoleh satu label, normal atau ugal-ugalan, sesuai perilaku pengemudi saat data direkam. Ukuran jendela menentukan panjang pola berkendara yang dilihat model dalam satu kali prediksi. Nilai $W$ dan $S$ penelitian ini ditetapkan pada BAB III.

\subsubsection{Pembagian Data dan Kebocoran Data}
Penelitian ini membagi dataset menjadi data latih, data validasi, dan data uji. Data latih dipakai untuk memperbarui bobot model. Data validasi dipakai untuk memantau pelatihan, memilih \textit{epoch} terbaik, dan menyetel \textit{hyperparameter}. Data uji dipakai satu kali pada akhir untuk mengukur kinerja model.

\textit{Stratified split} membagi dataset dengan menjaga perbandingan jumlah sampel antarkelas pada ketiga bagian. Pada data deret waktu, pembagian tidak boleh dilakukan secara acak per jendela. \textit{Sliding window} menghasilkan jendela yang saling tumpang tindih, sehingga dua jendela bertetangga hampir identik. Jika keduanya jatuh ke bagian yang berbeda, model sudah melihat isi data uji saat berlatih. Kondisi ini disebut kebocoran data (\textit{data leakage}) dan membuat hasil pengujian lebih tinggi dari kemampuan sebenarnya.

Kebocoran data dicegah dengan membagi data dalam satuan blok. Setiap segmen dipotong menjadi blok yang terdiri atas banyak cuplikan berurutan. Blok-blok inilah yang dibagi ke data latih, validasi, dan uji secara terpisah per kelas, lalu jendela dibuat di dalam setiap blok. Dengan cara ini, tidak ada jendela yang berbagi cuplikan dengan jendela di bagian lain. Pengujian yang paling ketat memakai sesi mengemudi utuh yang tidak dipakai saat pelatihan, karena blok-blok dari satu sesi tetap memiliki kondisi yang mirip, seperti rute, kendaraan, dan suasana hati pengemudi.

\subsection{Deep Learning dan Convolutional Neural Network}
\subsubsection{Jaringan Saraf Tiruan}
Jaringan saraf tiruan adalah model komputasi yang terdiri atas lapisan masukan, satu atau lebih lapisan tersembunyi, dan lapisan keluaran. Setiap neuron mengalikan setiap masukan dengan bobotnya, menjumlahkan hasilnya, menambahkan bias, lalu meneruskan hasilnya melalui fungsi aktivasi $\varphi$:
\begin{equation}
  z = w_{1}x_{1} + w_{2}x_{2} + \dots + w_{n}x_{n} + b, \qquad a = \varphi(z)
  \label{eq:neuron}
\end{equation}
dengan $x_{i}$ adalah masukan, $w_{i}$ adalah bobot, $b$ adalah bias, dan $a$ adalah keluaran neuron. Pelatihan mencari nilai $w_{i}$ dan $b$ yang membuat fungsi \textit{loss} $L$, yaitu ukuran kesalahan model, sekecil mungkin. Algoritma \textit{backpropagation} menghitung seberapa besar pengaruh setiap bobot terhadap \textit{loss}, yang disebut gradien $\partial L/\partial w$. Setelah itu, \textit{gradient descent} memperbarui bobot ke arah yang menurunkan \textit{loss}:
\begin{equation}
  w_{\mathrm{baru}} = w_{\mathrm{lama}} - \eta\,\frac{\partial L}{\partial w}
  \label{eq:gd}
\end{equation}
dengan $\eta$ adalah laju pembelajaran yang menentukan besar langkah pembaruan. Pada penelitian ini, lapisan keluaran terdiri atas satu neuron dengan aktivasi sigmoid
\begin{equation}
  \mathrm{sig}(z) = \frac{1}{1 + e^{-z}}
  \label{eq:sigmoid}
\end{equation}
Sigmoid mengubah bilangan berapa pun menjadi nilai antara 0 dan 1. Keluarannya, $\hat{y}$, dibaca sebagai probabilitas bahwa suatu jendela telemetri berasal dari perilaku ugal-ugalan.

\subsubsection{Arsitektur CNN (konvolusi, pooling, fully connected)}
\textit{Convolutional Neural Network} (CNN) menyusun tiga jenis lapisan secara berurutan. Lapisan konvolusi mengekstrak fitur lokal dari masukan. Lapisan \textit{pooling} meringkas fitur tersebut. Lapisan \textit{fully connected} menggabungkan fitur dan menghasilkan keputusan klasifikasi.

Lapisan konvolusi menggeser sebuah filter kecil di sepanjang masukan. Pada setiap posisi, filter menghitung jumlah berbobot dari nilai-nilai yang tercakup, sehingga hasilnya besar ketika pola pada data mirip dengan pola pada filter. Lapisan \textit{max pooling} berukuran 2 mengambil nilai terbesar dari setiap dua langkah berurutan. Operasi ini memperpendek deret menjadi setengahnya, misalnya dari 30 langkah menjadi 15 langkah, dan tetap menyimpan respons terkuat dari setiap pola. Lapisan \textit{fully connected} bekerja seperti neuron pada Persamaan~\ref{eq:neuron}, dan lapisan terakhirnya memakai sigmoid untuk memutuskan apakah jendela termasuk normal atau ugal-ugalan.

\subsubsection{CNN 1 Dimensi untuk Data Deret Waktu}
CNN satu dimensi (1D-CNN) menggeser filter hanya sepanjang sumbu waktu. Masukannya berupa jendela berukuran $W \times F$, yaitu $W$ langkah waktu dan $F$ fitur. Untuk satu fitur dan filter berukuran 3, keluaran pada langkah ke-$t$ adalah
\begin{equation}
  y_{t} = \varphi\big(w_{1}\,x_{t-2} + w_{2}\,x_{t-1} + w_{3}\,x_{t} + b\big)
  \label{eq:konv1d}
\end{equation}
dengan $x_{t}$ adalah nilai fitur pada langkah ke-$t$, $w_{1}$ sampai $w_{3}$ adalah bobot filter, dan $b$ adalah bias. Jika masukan memiliki 11 fitur, filter memiliki tiga bobot untuk setiap fitur, dan hasil perkalian dari seluruh fitur dijumlahkan menjadi satu nilai $y_{t}$. Filter yang sama dipakai di seluruh langkah waktu, sehingga jumlah parameternya kecil dan filter dapat mengenali pola yang sama di posisi mana pun dalam jendela.

Persamaan~\ref{eq:konv1d} hanya memakai cuplikan saat ini dan dua cuplikan sebelumnya. Bentuk ini disebut konvolusi kausal (\textit{causal convolution}). Agar langkah pertama tetap memiliki keluaran, dua nilai nol ditambahkan di awal deret (\textit{padding}), sehingga panjang keluaran sama dengan panjang masukan. Sifat kausal sesuai untuk data yang diproses secara \textit{real-time}, karena model tidak memakai informasi dari masa depan.

Pada data telemetri, satu filter dapat mempelajari pola seperti kenaikan akselerasi yang cepat disertai kenaikan RPM, atau penurunan kecepatan yang tajam setelah pedal gas dilepas. Pola-pola ini mencerminkan akselerasi agresif dan pengereman mendadak, sehingga 1D-CNN sesuai untuk mengekstrak pola lokal dari jendela telemetri.

\subsubsection{Recurrent Neural Network dan Long Short-Term Memory}
Konvolusi hanya melihat pola sepanjang ukuran filternya. Perilaku mengemudi juga memiliki ketergantungan yang lebih panjang, misalnya rangkaian akselerasi keras, pengereman, lalu belokan tajam dalam beberapa detik. \textit{Recurrent Neural Network} (RNN) memproses deret langkah demi langkah dan membawa keadaan tersembunyi $\mathbf{h}_{t}$ dari langkah sebelumnya, sehingga dapat merangkum urutan kejadian. RNN biasa sulit mempelajari ketergantungan jangka panjang karena gradiennya mengecil ketika dirambatkan balik melalui banyak langkah waktu.

\textit{Long Short-Term Memory} (LSTM) yang diusulkan Hochreiter dan Schmidhuber mengatasi masalah tersebut dengan sel memori $\mathbf{c}_{t}$ yang diatur oleh tiga gerbang. Gerbang lupa $\mathbf{f}_{t}$ menentukan bagian memori lama yang dipertahankan. Gerbang masukan $\mathbf{i}_{t}$ menentukan informasi baru $\tilde{\mathbf{c}}_{t}$ yang disimpan. Gerbang keluaran $\mathbf{o}_{t}$ menentukan bagian memori yang dikeluarkan. Setiap gerbang bernilai antara 0 dan 1 karena memakai sigmoid. Memori dan keluaran LSTM diperbarui dengan
\begin{equation}
  \mathbf{c}_{t} = \mathbf{f}_{t} \odot \mathbf{c}_{t-1} + \mathbf{i}_{t} \odot \tilde{\mathbf{c}}_{t}, \qquad
  \mathbf{h}_{t} = \mathbf{o}_{t} \odot \tanh(\mathbf{c}_{t})
  \label{eq:lstm}
\end{equation}
dengan $\odot$ adalah perkalian elemen demi elemen. Persamaan pertama berarti memori baru adalah sebagian memori lama ditambah sebagian informasi baru. Karena memori diperbarui dengan penjumlahan, gradien dapat mengalir melalui banyak langkah waktu tanpa cepat menghilang. Satu lapisan LSTM dengan $u$ unit dan $d$ masukan memiliki parameter sebanyak
\begin{equation}
  N_{\mathrm{LSTM}} = 4\,\big(u\,(d + u) + u\big)
  \label{eq:param-lstm}
\end{equation}
Angka 4 berasal dari tiga gerbang dan satu calon memori $\tilde{\mathbf{c}}_{t}$, yang masing-masing memiliki bobot sendiri.

\textit{Bidirectional} LSTM (BiLSTM) yang diusulkan Schuster dan Paliwal menjalankan dua LSTM pada deret yang sama, satu dari awal ke akhir dan satu dari akhir ke awal, lalu menggabungkan keadaan tersembunyi keduanya. Dengan demikian, representasi pada setiap langkah dalam satu jendela memuat konteks sebelum dan sesudahnya. Pada klasifikasi per jendela hal ini sah, karena seluruh isi jendela sudah tersedia saat prediksi dibuat.

\subsubsection{Arsitektur Hibrida CNN-LSTM}
Arsitektur hibrida CNN-LSTM menggabungkan kelebihan kedua jenis lapisan. Lapisan konvolusi mengekstrak pola lokal pendek dari setiap beberapa langkah waktu dan memperpendek deret melalui \textit{pooling}. Lapisan LSTM kemudian merangkum urutan pola lokal tersebut menjadi satu vektor yang mewakili seluruh jendela. Pada penelitian ini, CNN mengenali kejadian sesaat seperti lonjakan pedal atau akselerasi, sedangkan BiLSTM dan LSTM mengenali urutan kejadian-kejadian itu dalam jendela tiga detik. Jia dkk. memakai gabungan LSTM dan CNN untuk mengenali perilaku mengemudi abnormal dari data kendaraan nyata. Cura dkk. menunjukkan bahwa 1D-CNN dapat mengenali perilaku agresif dari data CAN bus.

\subsubsection{Normalisasi Lapisan}
\textit{Batch normalization} dari Ioffe dan Szegedy menormalkan keluaran suatu lapisan berdasarkan rata-rata dan varians dalam satu \textit{mini-batch}, lalu menskalakannya kembali dengan dua parameter yang dipelajari. Teknik ini menstabilkan sebaran nilai yang masuk ke lapisan berikutnya sehingga pelatihan lebih cepat dan tidak mudah divergen. \textit{Layer normalization} dari Ba dkk. melakukan normalisasi serupa, tetapi menghitung rata-rata dan varians dari seluruh unit pada satu sampel. Karena tidak bergantung pada ukuran \textit{batch}, \textit{layer normalization} cocok untuk lapisan rekuren.

\subsection{Fungsi Aktivasi}
\subsubsection{ReLU dan Masalah Dying ReLU}
\textit{Rectified Linear Unit} (ReLU) meneruskan masukan positif apa adanya dan mengubah masukan negatif menjadi nol:
\begin{equation}
  f(x) = \max(0, x)
  \label{eq:relu}
\end{equation}
Sebagai contoh, $f(2) = 2$, sedangkan $f(-3) = 0$. Untuk masukan negatif, keluaran ReLU tetap nol berapa pun nilai masukannya, sehingga gradiennya juga nol. Jika masukan suatu neuron negatif untuk semua data, bobot neuron tersebut tidak lagi menerima gradien, berhenti diperbarui, dan neuron tidak ikut belajar. Kondisi ini disebut \textit{dying ReLU}.

ReLU tidak bekerja langsung pada data masukan, tetapi pada jumlah berbobot $z$ di dalam setiap neuron sesuai Persamaan~\ref{eq:neuron}. Lapisan konvolusi pertama tetap menerima nilai akselerasi negatif apa adanya, lalu bobot filter dapat mengubah tanda nilai tersebut. Risiko \textit{dying ReLU} muncul ketika $z$ suatu neuron bernilai negatif untuk hampir seluruh data latih. Data telemetri yang telah dinormalisasi dengan \textit{Z-score} tersebar di sekitar nol, sehingga banyak neuron menerima $z$ negatif pada sebagian sampel. Pada model penelitian ini, filter konvolusi yang mati berhenti mendeteksi polanya, misalnya lonjakan akselerasi, sehingga kemampuan model membedakan perilaku mengemudi berkurang.

\subsubsection{Parametric ReLU (PReLU)}
\textit{Parametric ReLU} (PReLU) yang diusulkan He dkk. mempertahankan kemiringan pada bagian negatif dan menjadikan kemiringan itu parameter yang dipelajari:
\begin{equation}
  f(x) =
  \begin{cases}
    x, & x > 0 \\
    a\,x, & x \le 0
  \end{cases}
  \label{eq:prelu}
\end{equation}
Sebagai contoh, dengan $a = 0{,}25$, masukan $-4$ menghasilkan keluaran $-1$. Koefisien $a$ diperbarui bersama bobot jaringan melalui \textit{backpropagation}, dan nilainya hanya dipengaruhi oleh masukan negatif. Model dapat memakai satu nilai $a$ untuk seluruh lapisan atau satu nilai $a$ untuk setiap kanal. Karena keluaran untuk masukan negatif tidak lagi selalu nol, neuron tetap menerima sinyal pembelajaran walaupun masukannya negatif. Penelitian ini memakai PReLU agar model menyesuaikan perlakuan terhadap nilai negatif dengan karakter data telemetri, dan agar filter tidak mati seperti pada ReLU.

Pada model penelitian ini, PReLU dipakai pada lapisan konvolusi dan lapisan \textit{fully connected} tersembunyi. Lapisan LSTM tetap memakai sigmoid dan $\tanh$ di dalam gerbangnya sesuai Persamaan~\ref{eq:lstm}. PReLU hanya menambah satu parameter per kanal, sehingga tambahan biaya komputasinya kecil dibandingkan jumlah seluruh bobot model.


\subsection{Pelatihan dan Optimisasi Model}
\subsubsection{Fungsi Loss (Binary Cross-Entropy)}
Fungsi \textit{loss} mengukur seberapa jauh prediksi model dari label sebenarnya. Selama pelatihan, model mengubah bobotnya untuk menurunkan nilai \textit{loss}. Klasifikasi perilaku normal dan ugal-ugalan merupakan masalah dua kelas, sehingga penelitian ini memakai \textit{binary cross-entropy}:
\begin{equation}
  L = -\frac{1}{N} \sum_{i=1}^{N} \Big[ y_{i}\,\log \hat{y}_{i} + (1 - y_{i})\,\log (1 - \hat{y}_{i}) \Big]
  \label{eq:bce}
\end{equation}
dengan $N$ adalah jumlah sampel, $y_{i}$ adalah label sebenarnya (1 untuk ugal-ugalan, 0 untuk normal), dan $\hat{y}_{i}$ adalah probabilitas ugal-ugalan dari sigmoid. Untuk satu sampel ugal-ugalan, hanya suku pertama yang tersisa, sehingga $L = -\log \hat{y}$. Jika model memberi $\hat{y} = 0{,}9$, nilai \textit{loss}-nya sekitar 0,11. Jika model salah dan memberi $\hat{y} = 0{,}1$, nilai \textit{loss}-nya sekitar 2,30. Kesalahan dengan keyakinan tinggi mendapat \textit{loss} terbesar, sehingga mendapat koreksi terbesar saat pembaruan bobot.

\subsubsection{Optimizer (Adam) dan Hyperparameter}
\textit{Optimizer} adalah algoritma yang memperbarui bobot jaringan berdasarkan gradien \textit{loss}. \textit{Adaptive Moment Estimation} (Adam) dari Kingma dan Ba menyimpan rata-rata gradien dari langkah-langkah sebelumnya beserta rata-rata kuadratnya. Dari kedua nilai itu, Adam menyesuaikan besar langkah pembaruan untuk setiap parameter secara terpisah. Parameter dengan gradien yang sering berganti arah memperoleh langkah yang lebih kecil, sedangkan parameter dengan gradien yang konsisten memperoleh langkah yang mendekati laju pembelajaran. Mekanisme ini menjaga pelatihan tetap stabil tanpa penyetelan laju pembelajaran yang rumit. Laju pembelajaran juga dapat diturunkan bertahap ketika \textit{loss} validasi berhenti membaik, sehingga bobot mendekati titik optimum dengan langkah yang lebih halus.

\textit{Hyperparameter} adalah nilai yang peneliti tetapkan sebelum pelatihan dan tidak dipelajari dari data. Pada penelitian ini, \textit{hyperparameter} meliputi laju pembelajaran (\textit{learning rate}), \textit{batch size}, jumlah \textit{epoch}, ukuran dan pergeseran jendela, jumlah filter dan unit LSTM, serta laju \textit{dropout}. Penelitian ini memilih nilai-nilai tersebut berdasarkan kinerja pada data validasi, sehingga data uji tetap bebas dari proses penyetelan.

\subsubsection{Overfitting, Dropout, dan Early Stopping}
\textit{Overfitting} terjadi ketika model menghafal data latih sehingga kinerjanya turun pada data baru. Tandanya, \textit{loss} latih terus turun sementara \textit{loss} validasi naik. Risiko ini besar pada penelitian dengan dataset yang direkam dari sesi mengemudi yang terbatas.

\textit{Dropout} dari Srivastava dkk. menonaktifkan sebagian neuron secara acak pada setiap iterasi pelatihan. Dengan laju \textit{dropout} $p = 0{,}3$, sekitar 30\% neuron pada lapisan tersebut dinonaktifkan pada setiap iterasi, dan keluaran neuron yang tersisa diperbesar agar total sinyalnya tetap sama. Teknik ini memaksa jaringan tidak bergantung pada satu neuron saja. \textit{Dropout} hanya aktif saat pelatihan dan tidak dipakai saat inferensi. Karena itu, \textit{loss} latih dapat lebih tinggi daripada \textit{loss} validasi tanpa menandakan kesalahan.

\textit{Early stopping} menghentikan pelatihan ketika metrik validasi tidak membaik selama $n_{p}$ \textit{epoch} berturut-turut. Nilai $n_{p}$ disebut \textit{patience}. Setelah pelatihan berhenti, bobot dari \textit{epoch} dengan metrik validasi terbaik dikembalikan sebagai bobot model.

\subsubsection{Ketidakseimbangan Kelas (Class Imbalance)}
Ketidakseimbangan kelas terjadi ketika jumlah sampel antarkelas berbeda jauh. Rasio ketidakseimbangan dihitung dengan
\begin{equation}
  \mathrm{IR} = \frac{n_{\mathrm{mayoritas}}}{n_{\mathrm{minoritas}}}
  \label{eq:ir}
\end{equation}
dengan $n_{\mathrm{mayoritas}}$ dan $n_{\mathrm{minoritas}}$ adalah jumlah sampel kelas terbanyak dan kelas tersedikit. Model yang dilatih pada data seperti ini condong memprediksi kelas mayoritas, karena kelas tersebut paling banyak menyumbang \textit{loss}. Kecondongan ke kelas ugal-ugalan membuat banyak pengemudi normal ditandai berbahaya. Kecondongan ke kelas normal membuat pengemudi ugal-ugalan lolos tanpa terdeteksi.

Pembobotan kelas mengurangi masalah ini dengan memberi bobot lebih besar pada kelas minoritas dalam \textit{loss}. Bobot kelas $c$ dihitung dengan
\begin{equation}
  w_{c} = \frac{N}{C\,n_{c}}
  \label{eq:bobot-kelas}
\end{equation}
dengan $N$ adalah jumlah seluruh sampel, $C$ adalah jumlah kelas, dan $n_{c}$ adalah jumlah sampel kelas $c$. Sebagai contoh, dataset berisi 600 jendela normal dan 400 jendela ugal-ugalan memiliki IR 1,5. Bobot kelas normal adalah $1000/(2 \times 600) \approx 0{,}83$ dan bobot kelas ugal-ugalan adalah $1000/(2 \times 400) = 1{,}25$. Metode lain adalah \textit{oversampling} kelas minoritas, \textit{undersampling} kelas mayoritas, dan augmentasi data. Karena akurasi dapat tinggi walaupun kelas minoritas gagal dikenali, penelitian ini menilai model dengan presisi, \textit{recall}, dan \textit{F1-score} per kelas, bukan hanya akurasi.

\subsection{Evaluasi Kinerja Model}
\subsubsection{Confusion Matrix}
\textit{Confusion matrix} adalah tabel yang membandingkan prediksi model dengan label sebenarnya. Penelitian ini menetapkan ugal-ugalan sebagai kelas positif dan normal sebagai kelas negatif. Tabel~\ref{tab:confusion} menunjukkan empat nilai yang dihasilkan.

\begin{table}[H]
  \centering
  \caption{\textit{Confusion matrix} untuk klasifikasi perilaku mengemudi}
  \label{tab:confusion}
  \begin{tabular}{|l|c|c|}
    \hline
     & \textbf{Prediksi ugal-ugalan} & \textbf{Prediksi normal} \\
    \hline
    \textbf{Aktual ugal-ugalan} & TP & FN \\
    \hline
    \textbf{Aktual normal} & FP & TN \\
    \hline
  \end{tabular}
\end{table}

\textit{True positive} (TP) adalah jendela ugal-ugalan yang diprediksi ugal-ugalan. \textit{True negative} (TN) adalah jendela normal yang diprediksi normal. \textit{False positive} (FP) adalah jendela normal yang diprediksi ugal-ugalan, atau alarm palsu. \textit{False negative} (FN) adalah jendela ugal-ugalan yang diprediksi normal, atau perilaku berbahaya yang tidak terdeteksi.

\subsubsection{Akurasi, Presisi, Recall, dan F1-score}
Penelitian ini menghitung empat metrik dari nilai pada \textit{confusion matrix}. Akurasi adalah proporsi prediksi yang benar dari seluruh sampel:
\begin{equation}
  \mathrm{Akurasi} = \frac{TP + TN}{TP + TN + FP + FN}
  \label{eq:akurasi}
\end{equation}
Presisi adalah proporsi prediksi ugal-ugalan yang memang ugal-ugalan:
\begin{equation}
  \mathrm{Presisi} = \frac{TP}{TP + FP}
  \label{eq:presisi}
\end{equation}
\textit{Recall} adalah proporsi perilaku ugal-ugalan yang berhasil terdeteksi:
\begin{equation}
  \mathrm{Recall} = \frac{TP}{TP + FN}
  \label{eq:recall}
\end{equation}
\textit{F1-score} adalah rata-rata harmonik presisi dan \textit{recall}:
\begin{equation}
  F_{1} = \frac{2 \times \mathrm{Presisi} \times \mathrm{Recall}}{\mathrm{Presisi} + \mathrm{Recall}}
  \label{eq:f1}
\end{equation}

Presisi rendah berarti banyak pengemudi normal ditandai ugal-ugalan. \textit{Recall} rendah berarti banyak perilaku ugal-ugalan lolos. Pada data yang tidak seimbang, akurasi saja dapat menyesatkan, sehingga penelitian ini memakai \textit{F1-score} untuk menilai keseimbangan antara presisi dan \textit{recall}.

\subsubsection{Kurva ROC dan AUC}
Keluaran model berupa probabilitas $\hat{y}$, sehingga nilai presisi dan \textit{recall} bergantung pada ambang keputusan yang dipilih. Menurut Fawcett, kurva \textit{Receiver Operating Characteristic} (ROC) menggambarkan \textit{true positive rate} terhadap \textit{false positive rate} untuk seluruh kemungkinan ambang. Luas daerah di bawah kurva tersebut (\textit{Area Under the Curve}, AUC) dapat diartikan sebagai peluang model memberi skor lebih tinggi pada satu sampel positif acak daripada pada satu sampel negatif acak. Nilai AUC 0,5 setara dengan tebakan acak, dan nilai 1 berarti kedua kelas terpisah sempurna. AUC juga dapat dihitung untuk satu fitur saja, dengan memakai nilai fitur itu sebagai skor, untuk mengukur kemampuan fitur tersebut membedakan kedua kelas tanpa bantuan fitur lain.

\subsubsection{Latensi Inferensi untuk Sistem Real-time}
Latensi inferensi adalah waktu yang dibutuhkan model sejak menerima satu jendela masukan sampai menghasilkan prediksi. Pada sistem \textit{real-time}, hasil klasifikasi harus keluar sebelum prediksi berikutnya dijadwalkan. Jika waktu pemrosesan lebih lama dari selang tersebut, prediksi tertinggal dari kondisi mengemudi yang sedang berlangsung dan antrean data terus bertambah. Latensi total mencakup waktu prapemrosesan, yaitu normalisasi dan pembentukan jendela, ditambah waktu model menghitung prediksi.

Penelitian ini mengukur latensi dengan menjalankan inferensi berulang kali lalu mengambil rata-ratanya, karena waktu satu kali eksekusi dapat dipengaruhi beban sistem operasi pada saat itu. Hasil pengukuran menunjukkan apakah model hibrida CNN-BiLSTM dengan PReLU cukup ringan untuk dipakai pada sistem pemantau perilaku mengemudi.

  \cleardoublepage

  % Konten metodologi
  % !TEX root = ../main.tex
\chapter{METODOLOGI}

\section{Metode yang digunakan}
Penelitian ini menggunakan metode eksperimental. Pendekatan eksperimental dipilih untuk merancang, melatih, mengevaluasi, dan memperbaiki model klasifikasi perilaku pengemudi berbasis data telemetri kendaraan. Data latih dikumpulkan dari sesi mengemudi pada simulator BeamNG.drive dengan dua skenario, yaitu mengemudi normal dan mengemudi ugal-ugalan. Label setiap data ditentukan oleh skenario yang sedang dijalankan saat data direkam. Model yang telah dilatih kemudian diuji pada dua lingkungan, yaitu simulator dan mobil nyata yang datanya dibaca melalui adaptor OBD-II ELM327.

Variabel bebas pada eksperimen ini adalah kombinasi fitur masukan yang diuji melalui ablasi fitur. Variabel terikatnya adalah kinerja klasifikasi dan latensi inferensi. Variabel kontrolnya adalah arsitektur model dengan aktivasi PReLU, data latih, \textit{hyperparameter}, \textit{seed} acak, kendaraan, mode transmisi, serta kondisi lingkungan simulasi. Kinerja diukur secara kuantitatif dengan akurasi, presisi, \textit{recall}, \textit{F1-score}, AUC, dan latensi inferensi pada data uji, pengujian \textit{real-time} di simulator, dan pengujian pada mobil.

\section{Bahan dan peralatan yang digunakan}
Untuk merealisasikan sistem deteksi gaya berkendara berbahaya (ugal-ugalan) yang diusulkan, dibutuhkan komponen perangkat keras serta perangkat lunak sebagai berikut:

\subsection{Perangkat Keras (\textit{Hardware})}
\begin{enumerate}[label=\arabic*)]
  \item \textbf{Komputer Simulator}: PC spesifikasi tinggi yang menjalankan simulator BeamNG.drive, program perekam telemetri, proses pelatihan model, dan program klasifikasi \textit{real-time}.
  \item \textbf{\textit{Driving Simulator}}: Rangkaian set kemudi fisik (\textit{steering wheel}) dan pedal (\textit{throttle}, \textit{brake}) yang berfungsi sebagai perangkat masukan mekanis bagi pengemudi untuk mengendalikan kendaraan di dalam simulator secara realistis.
  \item \textbf{Adaptor ELM327 WiFi}: Adaptor yang dipasang pada port OBD-II mobil uji dan mengirim data kendaraan ke laptop melalui jaringan WiFi.
  \item \textbf{Mobil Uji}: Satu mobil penumpang dengan port OBD-II dan jenis transmisi yang sama dengan kendaraan di simulator.
  \item \textbf{Laptop}: Laptop yang menjalankan program klasifikasi di dalam mobil saat pengujian.
\end{enumerate}

\subsection{Perangkat Lunak (\textit{Software})}
\begin{enumerate}[label=\arabic*), nosep, leftmargin=1.5em]
  \item \textbf{Sistem Operasi}: Windows 11 untuk menjalankan BeamNG.drive dan program Python secara bersamaan.
  \item \textbf{Perangkat Lunak Simulasi}: BeamNG.drive versi 0.39 sebagai simulator \textit{high-fidelity} dengan mesin fisika \textit{soft-body} yang menghasilkan dinamika kendaraan dan data telemetri.
  \item \textbf{Protokol Telemetri Simulator}: OutGauge dan MotionSim, yaitu protokol UDP bawaan BeamNG.drive yang mengirim data kendaraan ke program Python melalui alamat \textit{localhost}.
  \item \textbf{Koneksi OBD-II}: Koneksi TCP dari program Python ke adaptor ELM327 melalui WiFi untuk mengirim permintaan PID dan menerima jawabannya.
  \item \textbf{Bahasa Pemrograman}: Python 3.11.
  \item \textbf{Pustaka dan Kerangka Kerja (\textit{Library \& Framework})}:
  \begin{enumerate}[label=\alph*), nosep, leftmargin=1.5em]
    \item \textbf{TensorFlow 2.21 dan Keras 3}: Merancang arsitektur hibrida CNN-BiLSTM dengan fungsi aktivasi PReLU, melatih, dan menjalankan inferensi model.
    \item \textbf{scikit-learn}: Normalisasi \textit{Z-score}, pengodean label, serta perhitungan \textit{confusion matrix} dan laporan klasifikasi.
    \item \textbf{Pandas dan NumPy}: Pembersihan data, segmentasi, \textit{resampling}, perhitungan akselerasi, dan pembentukan \textit{sliding window}.
    \item \textbf{Matplotlib dan Seaborn}: Visualisasi kurva pelatihan dan \textit{confusion matrix}.
  \end{enumerate}
\end{enumerate}

\section{Urutan Pelaksanaan Penelitian}
Pelaksanaan penelitian ini dibagi menjadi beberapa tahapan yang saling berurutan, mulai dari pengkajian penelitian terdahulu hingga evaluasi model. Alur kerja penelitian diilustrasikan pada Gambar~\ref{fig:urutan} dan dijelaskan pada setiap subbab berikut.
\begin{figure}[H]
  \centering
  \begin{tikzpicture}[
      node distance=3mm,
      langkah/.style={draw, rounded corners=3pt, minimum width=7.4cm, minimum height=6.2mm, align=center, font=\small},
      putusan/.style={draw, diamond, aspect=2.6, inner sep=1pt, align=center, font=\small},
      titik/.style={circle, fill=black, minimum size=3.5mm, inner sep=0pt},
      panah/.style={-{Stealth[length=2mm]}}]
    \node[titik] (mulai) {};
    \node[langkah, below=of mulai] (s1) {1. Pengkajian penelitian terdahulu};
    \node[langkah, below=of s1] (s2) {2. Penentuan parameter yang diambil};
    \node[langkah, below=of s2] (s3) {3. Pengambilan dataset telemetri di simulator};
    \node[langkah, below=of s3] (s4) {4. Ekstraksi fitur dan \textit{preprocessing}};
    \node[langkah, below=of s4] (s5) {5. Pembuatan algoritma CNN-BiLSTM dengan PReLU};
    \node[langkah, below=of s5] (s6) {6. \textit{Training} dan optimisasi model};
    \node[langkah, below=of s6] (s7) {7. \textit{Testing} pada BeamNG secara \textit{real-time}};
    \node[langkah, below=of s7] (s8) {8. \textit{Testing} pada mobil melalui OBD-II};
    \node[langkah, below=of s8] (s9) {9. Evaluasi model};
    \node[putusan, below=of s9] (cek) {Memenuhi\\kriteria?};
    \node[langkah, minimum width=3.2cm, below=of cek] (final) {Model final};
    \node[langkah, minimum width=3.2cm, right=12mm of cek] (perbaikan) {Perbaikan sistem};
    \node[titik, below=of final] (selesai) {};
    \node[circle, draw, minimum size=5mm, inner sep=0pt] at (selesai) {};
    \foreach \a/\b in {mulai/s1, s1/s2, s2/s3, s3/s4, s4/s5, s5/s6, s6/s7, s7/s8, s8/s9, s9/cek, final/selesai}
      \draw[panah] (\a) -- (\b);
    \draw[panah] (cek) -- node[right, font=\small] {Ya} (final);
    \draw[panah] (cek) -- node[above, font=\small] {Tidak} (perbaikan);
    \draw[panah] (perbaikan) |- (s6);
  \end{tikzpicture}
  \caption{\emph{Flowchart} penelitian yang akan dilakukan}
  \label{fig:urutan}
\end{figure}

\subsection{Pengkajian Penelitian Terdahulu}
Tahap pertama penelitian ini adalah mengkaji penelitian terdahulu sebagai landasan dan acuan arah penelitian. Pada tahap ini, peneliti mengumpulkan parameter yang dipakai untuk mengklasifikasi gaya berkendara. Parameter yang diperoleh adalah kecepatan, akselerasi dan deselerasi, akselerasi lateral, putaran mesin (RPM), posisi \textit{throttle}, serta kedalaman pedal rem. Muhamad Fadzil dkk. memakai putaran mesin, kecepatan, posisi \textit{throttle}, dan beban mesin dari OBD-II, sedangkan Rahayu memakai putaran mesin, kecepatan, beban mesin, dan suhu mesin. Akselerasi lateral memerlukan sensor gerak, seperti pada penelitian Carmona dkk. dan Singh dan Sharma.

\subsection{Penentuan Parameter dalam Penelitian}
Pada tahap ini ditentukan parameter yang menjadi masukan model. Parameter dipilih berdasarkan dua pertimbangan, yaitu relevansinya dengan indikator perilaku ugal-ugalan pada BAB II dan ketersediaannya di simulator dan di mobil uji. Model yang dilatih di simulator harus menerima fitur yang sama ketika diuji di mobil. Oleh karena itu, penelitian ini memakai empat fitur inti pada Tabel~\ref{tab:fitur}.

\begin{table}[H]
  \centering
  \caption{Fitur inti masukan model}
  \label{tab:fitur}
  \small
  \begin{tabular}{|c|l|l|l|c|}
    \hline
    \textbf{No.} & \textbf{Fitur} & \textbf{Simulator} & \textbf{Mobil (OBD-II)} & \textbf{Satuan} \\
    \hline
    1 & Kecepatan & OutGauge (\textit{speed}) & mode 01, PID 0D & km/jam \\
    \hline
    2 & Putaran mesin & OutGauge (\textit{rpm}) & mode 01, PID 0C & rpm \\
    \hline
    3 & Posisi \textit{throttle} & OutGauge (\textit{throttle}) & mode 01, PID 11 & \% \\
    \hline
    4 & Akselerasi longitudinal & selisih kecepatan & selisih kecepatan & m/s$^{2}$ \\
    \hline
  \end{tabular}
\end{table}

Akselerasi longitudinal dihitung dari selisih kecepatan dengan cara yang sama pada simulator dan mobil, lalu diratakan dengan Persamaan~\ref{eq:moving-average}. Kecepatan dari simulator dibulatkan ke km/jam bulat sebelum akselerasi dihitung, agar resolusinya sama dengan kecepatan dari OBD-II. Akselerasi dari MotionSim tidak menjadi masukan model. Nilai tersebut dipakai untuk memeriksa apakah akselerasi hasil selisih kecepatan mengikuti akselerasi yang dihitung mesin fisika simulator. Pengujian awal terhadap paket MotionSim menunjukkan bahwa akselerasi dinyatakan dalam sistem koordinat kendaraan. Komponen sumbu Y bernilai positif saat pedal gas ditekan dan negatif saat pengereman.

Sebelum pengambilan data mobil, program meminta PID 00, 20, dan 40 pada mode 01 untuk mengetahui PID yang didukung mobil uji. Jika mobil mendukung PID posisi pedal akselerator (PID 49), PID tersebut menggantikan PID 11 karena lebih setara dengan kedalaman pedal gas di simulator. Jika mobil juga menyediakan posisi gigi atau kedalaman pedal rem, besaran tersebut dapat ditambahkan sebagai fitur pada simulator dan mobil, lalu model dilatih ulang.

\subsection{Pengambilan Dataset Telemetri}
Data latih diekstrak dari simulator melalui protokol OutGauge (port 4444) dan MotionSim (port 4123) pada alamat \textit{localhost}. Program perekam Python membuka dua \textit{socket} UDP dan menyimpan paket terbaru dari setiap protokol. Paket MotionSim divalidasi dengan memeriksa penanda \texttt{"BNG1"} dan panjang paket 88 \textit{byte}. Pada setiap detak jam 10~Hz, program menggabungkan paket terbaru dari kedua protokol lalu menulis satu baris ke berkas \texttt{.csv}. Jika paket dari salah satu protokol lebih tua dari 0,5~s, baris tersebut dilewati agar tidak ada baris yang memuat data basi.

Data direkam oleh peneliti dan dua sampai tiga pengemudi lain. Setiap pengemudi menjalankan beberapa sesi untuk setiap kelas dengan kriteria yang sama, dan setiap sesi berdurasi paling lama 15 menit untuk menjaga fokus pengemudi. Pada sesi normal, pengemudi berkendara dengan kecepatan wajar, perpindahan gigi yang halus, serta akselerasi dan pengereman yang lembut. Pada sesi ugal-ugalan, pengemudi melakukan akselerasi dan pengereman mendadak, menahan gigi rendah, berbelok tajam, serta berpindah jalur berulang kali.

Sesi ugal-ugalan dapat memuat potongan yang tampak normal, misalnya saat kendaraan berbelok pelan di tikungan. Untuk mengurangi label yang keliru, pengemudi mempertahankan pola mengemudi agresif sepanjang sesi. Penyaringan kecepatan di bawah 15~km/jam juga membuang cuplikan saat kendaraan nyaris diam. Jendela yang salah diklasifikasi pada data uji diperiksa ulang untuk melihat apakah kesalahan berasal dari model atau dari label.

Kendaraan, mode transmisi, peta, cuaca, dan waktu dibuat sama untuk kedua kelas agar perbedaan data hanya berasal dari perilaku pengemudi. Rekaman dihentikan jika terjadi tabrakan, karena lonjakan akselerasi akibat benturan bukan bagian dari perilaku mengemudi. Setiap kelas ditargetkan memiliki tiga sampai lima sesi dengan rute yang berbeda. Data uji diambil dari sesi yang tidak dipakai saat pelatihan, dan bila memungkinkan dari pengemudi yang datanya tidak dipakai saat pelatihan. Hasil pendahuluan pada BAB IV masih berasal dari satu pengemudi, yaitu peneliti.

Rekaman setiap kelas disimpan dalam berkas \texttt{.csv} terpisah, yaitu \texttt{Normal.csv} untuk kelas normal dan \texttt{Ugal.csv} untuk kelas ugal-ugalan, lalu digabungkan pada tahap berikutnya.

\subsection{Ekstraksi Fitur dan \textit{Preprocessing}}
Data mentah yang terkumpul kemudian memasuki tahap ekstraksi fitur dan \textit{preprocessing} sebelum digunakan sebagai masukan model. Tahapan yang dilakukan meliputi:
\begin{enumerate}
    \item \textbf{Validasi skema dan pembersihan data:} Setiap berkas diperiksa kelengkapan kolomnya. Baris yang mengandung nilai kosong atau tidak hingga dihapus.
    \item \textbf{Penyaringan netral:} Baris dengan kecepatan di bawah 15 km/jam dibuang dari kedua kelas, karena kendaraan yang diam atau hampir diam tidak memperlihatkan gaya mengemudi. Tidak ada filter yang hanya diterapkan pada satu kelas. Pada percobaan awal, kelas normal disaring dengan batas RPM, kedalaman pedal gas, dan kecepatan. Filter tersebut dihapus karena membuat model dapat membedakan kelas hanya dengan mengenali batas filter, seperti yang dijelaskan pada BAB II.
    \item \textbf{Segmentasi:} Rekaman dipotong menjadi segmen-segmen yang barisnya berurutan. Segmen baru dimulai jika ada baris yang dibuang, waktu mundur, atau selisih waktu antarbaris melebihi 1~s. Setiap segmen hanya berisi satu kelas.
    \item \textbf{\textit{Resampling}:} Data simulator di dalam setiap segmen di-\textit{resample} dari 10~Hz ke laju cuplik mobil $f_{s}$ yang diukur pada Subbab~\ref{subsec:uji-mobil}.
    \item \textbf{Perhitungan akselerasi:} Kecepatan dibulatkan ke km/jam bulat, lalu akselerasi longitudinal dihitung dari selisih kecepatan dan diratakan dengan Persamaan~\ref{eq:moving-average}. Beberapa baris pertama setiap segmen tidak memiliki nilai rata-rata lengkap dan dibuang.
    \item \textbf{Pemotongan nilai ekstrem:} Nilai akselerasi longitudinal dipotong pada $\pm 15$~m/s$^{2}$.
    \item \textbf{Pembagian dataset per blok:} Setiap segmen dipotong menjadi blok berdurasi sekitar 20~s. Blok dari setiap kelas diacak lalu dibagi menjadi 70\% data latih, 15\% data validasi, dan 15\% data uji. Pembagian per blok dan per kelas mencegah kebocoran data antarjendela yang tumpang tindih sekaligus menjaga proporsi kelas.
    \item \textbf{Pembuatan \textit{sliding window}:} Jendela berdurasi 3~s dengan pergeseran 0,5~s dibuat di dalam setiap blok. Dalam cuplikan, ukuran jendela adalah $W = 3 f_{s}$ dan pergeserannya adalah $S = 0{,}5 f_{s}$ yang dibulatkan. Setiap jendela berbentuk matriks $W \times 4$ dan diberi label sesuai kelas bloknya.
    \item \textbf{Normalisasi data:} Rata-rata $\mu$ dan simpangan baku $\sigma$ setiap fitur dihitung dari data latih, lalu dipakai untuk menormalisasi data latih, validasi, uji, data \textit{real-time} simulator, dan data mobil dengan \textit{Z-score} (Persamaan~\ref{eq:zscore}).
\end{enumerate}
Daftar fitur, laju cuplik, ukuran jendela, ukuran perataan, dan batas pemotongan disimpan dalam satu berkas konfigurasi bersama parameter normalisasi. Program klasifikasi di simulator dan di mobil membaca berkas yang sama, sehingga urutan dan perlakuan fitur pada saat inferensi selalu sama dengan saat pelatihan.

\subsection{Pembuatan Algoritma CNN-BiLSTM dengan PReLU}
Tahap ini merancang model yang menerima satu jendela telemetri berukuran $W \times 4$ dan menghasilkan probabilitas bahwa jendela tersebut berasal dari perilaku ugal-ugalan. Model memakai arsitektur hibrida. Blok konvolusi mengekstrak pola lokal pendek, lalu lapisan BiLSTM dan LSTM merangkum urutan pola tersebut sepanjang jendela. Tabel~\ref{tab:arsitektur} merinci setiap lapisan. Sebagai contoh, jika $f_{s} = 5$~Hz, maka $W = 15$ dan keluaran \textit{max pooling} berisi 7 langkah.

\begin{table}[H]
  \centering
  \caption{Arsitektur model hibrida CNN-BiLSTM dengan PReLU}
  \label{tab:arsitektur}
  \small
  \begin{tabular}{|c|l|p{4.3cm}|c|r|}
    \hline
    \textbf{No.} & \textbf{Lapisan} & \textbf{Konfigurasi} & \textbf{Keluaran} & \textbf{Parameter} \\
    \hline
    1 & Masukan & jendela telemetri & $W \times 4$ & 0 \\
    \hline
    2 & Konvolusi 1D & 32 filter, kernel 3, \textit{padding} kausal & $W \times 32$ & 416 \\
    \hline
    3 & PReLU & satu $a$ per filter & $W \times 32$ & 32 \\
    \hline
    4 & \textit{Batch normalization} & -- & $W \times 32$ & 128 \\
    \hline
    5 & \textit{Max pooling} & ukuran 2 & $\lfloor W/2 \rfloor \times 32$ & 0 \\
    \hline
    6 & \textit{Dropout} & $p = 0{,}3$ & $\lfloor W/2 \rfloor \times 32$ & 0 \\
    \hline
    7 & BiLSTM & 32 unit per arah, keluaran per langkah & $\lfloor W/2 \rfloor \times 64$ & 16.640 \\
    \hline
    8 & \textit{Layer normalization} & -- & $\lfloor W/2 \rfloor \times 64$ & 128 \\
    \hline
    9 & LSTM & 16 unit, keluaran langkah terakhir & $16$ & 5.184 \\
    \hline
    10 & \textit{Layer normalization} & -- & $16$ & 32 \\
    \hline
    11 & \textit{Fully connected} & 16 neuron & $16$ & 272 \\
    \hline
    12 & PReLU & satu $a$ per neuron & $16$ & 16 \\
    \hline
    13 & \textit{Dropout} & $p = 0{,}3$ & $16$ & 0 \\
    \hline
    14 & \textit{Fully connected} + sigmoid & 1 neuron & $1$ & 17 \\
    \hline
    \multicolumn{4}{|r|}{\textbf{Total}} & \textbf{22.865} \\
    \hline
  \end{tabular}
\end{table}

Ukuran lapisan sengaja dibuat kecil karena jumlah jendela pada dataset berada pada orde ribuan, sehingga model yang besar mudah menghafal data latih. Konvolusi memakai \textit{padding} kausal agar keluaran setiap langkah hanya bergantung pada cuplikan saat ini dan sebelumnya. Kernel berukuran 3 dipilih karena pola yang membedakan perilaku ugal-ugalan, seperti lonjakan akselerasi dan penekanan pedal gas yang mendadak, terjadi dalam beberapa cuplikan berurutan. \textit{Max pooling} memperpendek deret menjadi setengahnya sebelum masuk ke lapisan rekuren. Lapisan BiLSTM dan LSTM memakai \textit{dropout} 0,2 pada masukannya, dan lapisan BiLSTM juga memakai \textit{recurrent dropout} 0,1.

Jumlah parameter lapisan konvolusi dengan kernel $K$, kanal masukan $F_{\mathrm{in}}$, dan filter $F_{\mathrm{out}}$ adalah
\begin{equation}
  N_{\mathrm{konv}} = (K\,F_{\mathrm{in}} + 1)\,F_{\mathrm{out}}
  \label{eq:met-param}
\end{equation}
Pada lapisan konvolusi dengan $K = 3$, $F_{\mathrm{in}} = 4$, dan $F_{\mathrm{out}} = 32$, jumlah parameternya adalah $(3 \times 4 + 1) \times 32 = 416$. Jumlah parameter BiLSTM dihitung dengan Persamaan~\ref{eq:param-lstm} untuk setiap arah, yaitu $2 \times 4\,(32\,(32 + 32) + 32) = 16.640$. Jumlah parameter tidak bergantung pada $W$, karena filter dan bobot LSTM yang sama dipakai di setiap langkah waktu. PReLU memakai satu koefisien $a$ untuk setiap filter atau neuron (Persamaan~\ref{eq:prelu}), sehingga seluruh PReLU hanya menambah 48 parameter, atau sekitar 0,2\% dari total parameter model.

Algoritma model bekerja dengan langkah berikut.
\begin{enumerate}
  \item Model menerima satu jendela yang sudah dinormalisasi.
  \item Lapisan konvolusi kausal mengekstrak pola pendek dari kombinasi keempat fitur, lalu PReLU, \textit{batch normalization}, \textit{max pooling}, dan \textit{dropout} memprosesnya.
  \item Lapisan BiLSTM membaca langkah-langkah hasil konvolusi dari dua arah dan menghasilkan representasi setiap langkah yang memuat konteks sebelum dan sesudahnya.
  \item Lapisan LSTM merangkum seluruh langkah menjadi satu vektor berukuran 16.
  \item Lapisan \textit{fully connected} dengan PReLU dan \textit{dropout} menggabungkan vektor tersebut menjadi representasi akhir.
  \item Neuron keluaran dengan sigmoid menghasilkan probabilitas $\hat{y}$ bahwa jendela tergolong ugal-ugalan.
\end{enumerate}

\subsection{\textit{Training} Model dan Optimisasi Model}
Tahap ini melatih model pada data latih dan menyetel \textit{hyperparameter} dengan data validasi. Data uji tidak dipakai pada tahap ini. Pelatihan memakai \textit{mini-batch}. Pada setiap iterasi, model menghitung probabilitas $\hat{y}$ untuk satu \textit{batch} jendela, menghitung \textit{loss}, lalu memperbarui bobot dengan \textit{optimizer}.

Fungsi \textit{loss} yang dipakai adalah \textit{binary cross-entropy} pada Persamaan~\ref{eq:bce} dengan pembobotan kelas. Setiap suku \textit{loss} dikalikan dengan bobot kelas sampel tersebut, yaitu $w_{1}$ untuk kelas ugal-ugalan dan $w_{0}$ untuk kelas normal. Kedua bobot dihitung dari data latih dengan Persamaan~\ref{eq:bobot-kelas}. Pembobotan ini menjaga agar model tidak condong ke salah satu kelas walaupun jumlah jendela kedua kelas sedikit berbeda.

\textit{Optimizer} yang dipakai adalah Adam. Tabel~\ref{tab:hyperparameter} memuat nilai yang dipakai pada percobaan awal dan rentang pencarian untuk tahap optimisasi.

\begin{table}[H]
  \centering
  \caption{\textit{Hyperparameter} pelatihan model}
  \label{tab:hyperparameter}
  \small
  \begin{tabular}{|l|c|c|}
    \hline
    \textbf{\textit{Hyperparameter}} & \textbf{Nilai awal} & \textbf{Rentang pencarian} \\
    \hline
    Laju pembelajaran awal $\eta$ & $10^{-3}$ & $\{10^{-3}, 10^{-4}\}$ \\
    \hline
    Penurunan $\eta$ (\textit{loss} validasi stagnan 8 \textit{epoch}) & $\times 0{,}5$, min. $10^{-6}$ & tetap \\
    \hline
    \textit{Batch size} & 32 & $\{32, 64\}$ \\
    \hline
    \textit{Dropout} lapisan konvolusi dan \textit{dense} $p$ & 0,3 & $\{0{,}3;\ 0{,}5\}$ \\
    \hline
    \textit{Dropout} LSTM / \textit{recurrent dropout} & 0,2 / 0,1 & tetap \\
    \hline
    $\beta_{1}$, $\beta_{2}$, $\epsilon$ (Adam) & 0,9; 0,999; $10^{-7}$ & tetap \\
    \hline
    Jumlah \textit{epoch} maksimum & 100 & tetap \\
    \hline
    \textit{Patience} \textit{early stopping} & 20 & $\{10, 20\}$ \\
    \hline
    Durasi jendela / pergeseran & 3~s / 0,5~s & durasi $\in \{2, 3, 5\}$~s \\
    \hline
    Ukuran perataan akselerasi $m$ & 3 cuplikan & $\{1, 3, 5\}$ \\
    \hline
    Durasi blok pembagian data & 20~s & tetap \\
    \hline
    \textit{Seed} acak & 42 & beberapa \textit{seed} \\
    \hline
  \end{tabular}
\end{table}

Pelatihan berhenti lebih awal jika AUC validasi tidak membaik selama 20 \textit{epoch}. Bobot dari \textit{epoch} dengan AUC validasi tertinggi disimpan sebagai bobot model. Laju pembelajaran diturunkan setengahnya jika \textit{loss} validasi tidak membaik selama 8 \textit{epoch}. Kurva \textit{loss}, akurasi, dan AUC data latih serta validasi terhadap \textit{epoch} dipantau untuk mendeteksi \textit{overfitting}. Jika \textit{loss} latih turun sementara \textit{loss} validasi naik, peneliti menaikkan $p$ atau mengurangi kapasitas model.

AUC validasi dapat mencapai nilai maksimum 1 sejak awal pelatihan, terutama jika data validasi berasal dari sesi yang sama dengan data latih. Pada kondisi itu, AUC tidak lagi dapat membedakan \textit{epoch} yang baik dari yang buruk. Oleh karena itu, setelah data dari beberapa sesi tersedia, pemilihan \textit{epoch} terbaik juga memakai \textit{loss} validasi.

Setiap konfigurasi dilatih ulang beberapa kali dengan nilai \textit{seed} acak yang berbeda. Hasilnya dilaporkan sebagai rata-rata dan simpangan baku, sehingga hasil tidak bergantung pada satu inisialisasi bobot yang kebetulan baik. Konfigurasi terbaik dipilih berdasarkan \textit{F1-score} data validasi. Bobot model, parameter normalisasi, pengodean label, dan berkas konfigurasi fitur disimpan untuk dipakai pada tahap pengujian.

Jika hasil evaluasi pada tahap berikutnya belum memenuhi kriteria, penelitian kembali ke tahap ini setelah peneliti memperbaiki sistem, sesuai alur pada Gambar~\ref{fig:urutan}.

\subsection{\textit{Testing} Simulasi pada BeamNG secara \textit{Real-Time}}
\label{subsec:uji-beamng}
Tahap ini menguji model terlatih pada aliran data telemetri langsung dari BeamNG. Pengujian ini berbeda dari pengujian pada data uji \textit{offline} karena model harus memproses data yang terus mengalir dan memberi hasil sebelum prediksi berikutnya dijadwalkan. Pengemudi mengemudikan kendaraan di BeamNG memakai set kemudi dan pedal, sedangkan program Python menjalankan klasifikasi.

Program klasifikasi bekerja dengan langkah berikut.
\begin{enumerate}
  \item BeamNG mengirim paket OutGauge melalui UDP ke program Python.
  \item Pada setiap detak jam dengan laju $f_{s}$, program mengambil kecepatan, RPM, dan posisi \textit{throttle} terbaru, membulatkan kecepatan, lalu menghitung akselerasi dengan Persamaan~\ref{eq:moving-average}.
  \item Program memotong nilai ekstrem dengan batas yang sama seperti saat pelatihan, lalu menyimpan baris tersebut pada \textit{buffer} berukuran $W$.
  \item Setiap ada cuplikan baru, dan paling cepat setiap 0,3~s, program menormalisasi isi \textit{buffer} memakai parameter dari data latih dan menghitung probabilitas $\hat{y}$.
  \item Program merata-ratakan probabilitas dalam 4,5 detik terakhir menjadi $\bar{y}$ untuk mengurangi keputusan yang berubah-ubah.
  \item Program menentukan kelas dengan histeresis, lalu menampilkan hasil pada terminal dan mencatat waktu, $\hat{y}$, $\bar{y}$, label sebenarnya, serta latensi ke berkas CSV.
\end{enumerate}
Keputusan kelas memakai dua ambang, yaitu ambang atas $\tau_{\mathrm{atas}} = 0{,}85$ dan ambang bawah $\tau_{\mathrm{bawah}} = 0{,}50$, dengan aturan berikut:
\begin{enumerate}
  \item Jika status saat ini normal, status berubah menjadi ugal-ugalan hanya ketika $\bar{y}$ melebihi 0,85.
  \item Jika status saat ini ugal-ugalan, status kembali menjadi normal hanya ketika $\bar{y}$ turun di bawah 0,50.
  \item Selain kedua kondisi tersebut, status tidak berubah.
\end{enumerate}
Aturan dua ambang ini disebut histeresis. Tanpa histeresis, status dapat berganti-ganti pada setiap prediksi ketika $\bar{y}$ berada di sekitar satu ambang. Kedua ambang dapat disesuaikan pada data validasi jika model masih condong ke satu kelas.

Pengujian memakai sesi mengemudi baru yang tidak termasuk dalam dataset pelatihan. Pengemudi menjalankan sesi normal dan sesi ugal-ugalan dengan kriteria yang sama seperti saat pengambilan data. Label sebenarnya mengikuti jenis sesi yang dijalankan.

Pengujian mengukur tiga hal. Pertama, kebenaran prediksi dihitung dengan mencocokkan kelas hasil prediksi dengan label sesi, lalu hasilnya disajikan dalam \textit{confusion matrix}. Kedua, latensi total setiap prediksi dihitung dari waktu \textit{buffer} diambil sampai prediksi keluar:
\begin{equation}
  T_{\mathrm{total}} = T_{\mathrm{pra}} + T_{\mathrm{inf}}
  \label{eq:met-latensi}
\end{equation}
dengan $T_{\mathrm{pra}}$ adalah waktu normalisasi dan pembentukan jendela, dan $T_{\mathrm{inf}}$ adalah waktu inferensi model. Sistem memenuhi syarat \textit{real-time} jika $T_{\mathrm{total}}$ tidak melebihi selang antara dua prediksi dan tidak lebih dari 0,3~s. Penelitian ini melaporkan nilai rata-rata dan nilai maksimum $T_{\mathrm{total}}$. Ketiga, pengujian mengukur waktu deteksi, yaitu selang antara awal sesi ugal-ugalan dan saat status pertama kali berubah menjadi ugal-ugalan. Waktu deteksi mencakup panjang jendela, perataan probabilitas, dan ambang histeresis, sehingga nilainya lebih besar daripada latensi komputasi.

\subsection{\textit{Testing} pada Mobil melalui OBD-II}
\label{subsec:uji-mobil}
Tahap ini menguji model pada mobil nyata. Adaptor ELM327 dipasang pada port OBD-II mobil uji. Laptop terhubung ke jaringan WiFi adaptor, lalu program Python membuka koneksi TCP ke adaptor tersebut. Program klasifikasi bekerja dengan langkah berikut.
\begin{enumerate}
  \item Program menginisialisasi ELM327 dan meminta PID 00, 20, dan 40 pada mode 01 untuk mengetahui PID yang didukung mobil uji.
  \item Program membaca PID 0D, 0C, dan 11 secara berulang selama beberapa menit untuk mengukur laju cuplik yang tercapai. Laju ini menjadi $f_{s}$ untuk \textit{resampling} data simulator dan untuk pelatihan model final.
  \item Saat pengujian, program meminta ketiga PID secara berurutan. Satu putaran permintaan menjadi satu cuplikan, dan waktu tiba setiap jawaban dicatat.
  \item Program menghitung akselerasi dengan Persamaan~\ref{eq:moving-average}, memotong nilai ekstrem, menormalisasi data memakai parameter dari data latih, lalu mengisi \textit{buffer} berukuran $W$.
  \item Prediksi, perataan probabilitas, dan keputusan histeresis memakai cara yang sama dengan pengujian di BeamNG.
\end{enumerate}

Pengujian dilakukan di area tertutup yang bebas dari kendaraan dan pejalan kaki lain, dengan izin pengelola area. Sesi normal berisi akselerasi dan pengereman yang lembut serta kecepatan yang stabil. Sesi agresif terkendali berisi akselerasi penuh dari kondisi berhenti, pengereman keras, serta akselerasi dan pengereman yang berulang, dengan batas kecepatan yang disepakati bersama pembimbing. Pengemudi memakai sabuk pengaman, dan orang kedua di kursi penumpang menjalankan laptop agar pengemudi tetap fokus. Label sebenarnya mengikuti jenis sesi yang dijalankan.

Kecepatan di area tertutup lebih rendah daripada kecepatan di simulator. Hasil pengujian di mobil juga menunjukkan seberapa jauh keputusan model bergantung pada kecepatan. Pengujian ini mengukur kebenaran prediksi, latensi, dan waktu deteksi dengan cara yang sama seperti pengujian di BeamNG.

\subsection{Evaluasi Model}
Evaluasi menilai model pada empat sumber data, yaitu data uji hasil pembagian per blok, sesi uji simulator yang tidak dipakai saat pelatihan, pengujian \textit{real-time} di BeamNG, dan pengujian di mobil. Evaluasi memakai \textit{confusion matrix} dengan ugal-ugalan sebagai kelas positif, lalu menghitung akurasi, presisi, \textit{recall}, dan \textit{F1-score} dengan Persamaan~\ref{eq:akurasi}, \ref{eq:presisi}, \ref{eq:recall}, dan \ref{eq:f1}, serta AUC. Presisi dan \textit{recall} dihitung untuk kedua kelas karena kecondongan model ke satu kelas hanya terlihat jika kedua kelas dinilai terpisah. \textit{Recall} kelas ugal-ugalan yang rendah menunjukkan banyak perilaku berbahaya yang lolos, sedangkan presisi kelas ugal-ugalan yang rendah menunjukkan banyak pengemudi normal yang salah ditandai.

Selain metrik kinerja, evaluasi memuat tiga analisis tambahan:
\begin{enumerate}
  \item \textbf{Ablasi fitur:} model dilatih ulang tanpa fitur tertentu, misalnya tanpa akselerasi atau tanpa posisi \textit{throttle}, untuk mengetahui kontribusi setiap fitur.
  \item \textbf{Analisis daya pisah fitur:} AUC setiap fitur tunggal dihitung untuk melihat fitur mana yang paling menentukan keputusan model, dan untuk memeriksa apakah model bergantung pada fitur yang dapat menjadi pengganggu.
  \item \textbf{Pemeriksaan sebaran data mobil:} sebaran setiap fitur pada data mobil diperiksa terhadap rentang data latih, untuk melihat apakah data mobil berada di luar rentang yang pernah dilihat model.
\end{enumerate}

Alur penelitian pada Gambar~\ref{fig:urutan} menetapkan bahwa model diterima sebagai model final jika memenuhi kriteria. Penelitian ini menetapkan lima kriteria pada Tabel~\ref{tab:kriteria}, dan model harus memenuhi kelimanya.

\begin{table}[H]
  \centering
  \caption{Kriteria penerimaan model}
  \label{tab:kriteria}
  \begin{tabular}{|c|p{8.5cm}|c|}
    \hline
    \textbf{No.} & \textbf{Kriteria} & \textbf{Target} \\
    \hline
    1 & Akurasi pada sesi uji simulator dan pengujian \textit{real-time} di BeamNG & $\ge 90\%$ \\
    \hline
    2 & \textit{F1-score} setiap kelas pada data simulator & $\ge 0{,}85$ \\
    \hline
    3 & Akurasi pada pengujian di mobil & $\ge 80\%$ \\
    \hline
    4 & \textit{F1-score} setiap kelas pada pengujian di mobil & $\ge 0{,}75$ \\
    \hline
    5 & Latensi total rata-rata & $\le 0{,}3$~s \\
    \hline
  \end{tabular}
\end{table}

Target pada mobil lebih rendah daripada target pada simulator karena model belum pernah melihat data mobil saat pelatihan. Jika seluruh kriteria terpenuhi, model disimpan sebagai model final. Jika salah satu kriteria belum terpenuhi, peneliti memperbaiki sistem sesuai kriteria yang gagal, lalu mengulang tahap \textit{training}. Perbaikan yang dapat dilakukan adalah:
\begin{enumerate}
  \item menambah sesi rekaman pada kelas yang \textit{recall}-nya rendah,
  \item mengubah bobot kelas atau ambang $\tau_{\mathrm{atas}}$ dan $\tau_{\mathrm{bawah}}$ jika model condong ke satu kelas,
  \item mengubah durasi jendela atau ukuran perataan akselerasi jika pola berkendara tidak tertangkap dengan baik,
  \item menambah fitur dari PID lain yang didukung mobil uji, atau menyesuaikan rentang kecepatan data latih, jika kinerja pada mobil jauh di bawah kinerja pada simulator,
  \item mengubah \textit{dropout}, jumlah filter, atau jumlah unit LSTM jika terjadi \textit{overfitting},
  \item memperkecil model jika latensi melebihi 0,3~s.
\end{enumerate}

  \cleardoublepage

  % Konten lainnya
  % !TEX root = ../main.tex
\chapter{HASIL YANG DIHARAPKAN}

\section{Hasil yang Diharapkan dari Penelitian}

Dari penelitian yang akan dilakukan, diharapkan diperoleh hasil sebagai berikut:
\begin{enumerate}
  \item Sistem akuisisi telemetri yang merekam data OutGauge dan MotionSim dari BeamNG.drive pada frekuensi tetap 10~Hz, beserta dataset berlabel normal dan ugal-ugalan dari beberapa sesi mengemudi.
  \item Program akuisisi OBD-II yang membaca kecepatan, RPM, dan posisi \textit{throttle} dari mobil melalui adaptor ELM327 WiFi, beserta hasil pengukuran laju cuplik yang tercapai.
  \item Model hibrida CNN-BiLSTM dengan aktivasi PReLU yang memenuhi kriteria pada Tabel~\ref{tab:kriteria}, yaitu akurasi minimal 90\% dan \textit{F1-score} minimal 0,85 untuk setiap kelas pada data simulator, serta akurasi minimal 80\% dan \textit{F1-score} minimal 0,75 untuk setiap kelas pada pengujian di mobil.
  \item Analisis kontribusi setiap fitur terhadap keputusan model.
  \item Program klasifikasi \textit{real-time} untuk simulator dan mobil yang menghasilkan prediksi dengan latensi rata-rata tidak lebih dari 0,3~s.
\end{enumerate}

\section{Hasil Pendahuluan}

Sampai saat ini telah dilakukan dua percobaan pelatihan di simulator dengan susunan lapisan pada Tabel~\ref{tab:arsitektur}, \textit{hyperparameter} awal pada Tabel~\ref{tab:hyperparameter}, dan pembagian data per blok. Kedua percobaan ini dilakukan sebelum rencana pengujian di mobil ditetapkan, sehingga masih memakai data 10~Hz dengan jendela 30 cuplikan dan fitur dari OutGauge serta MotionSim. Kedua percobaan berbeda pada data, fitur, dan penyaringan yang dipakai. Tabel~\ref{tab:kondisi-percobaan} merangkum perbedaan tersebut.

\begin{table}[H]
  \centering
  \caption{Kondisi kedua percobaan awal}
  \label{tab:kondisi-percobaan}
  \small
  \begin{tabular}{|p{3.6cm}|p{5cm}|p{5cm}|}
    \hline
    \textbf{Aspek} & \textbf{Percobaan 1} & \textbf{Percobaan 2} \\
    \hline
    Jumlah fitur & 8: kecepatan, RPM, gigi, pedal gas, pedal rem, akselerasi longitudinal dan lateral, laju yaw & 11: delapan fitur tersebut ditambah sudut kemudi hasil estimasi dan dua turunannya \\
    \hline
    Pembacaan MotionSim & Struktur paket belum sesuai; kolom laju yaw berisi posisi vertikal kendaraan & Sesuai dokumentasi dan telah diverifikasi \\
    \hline
    Filter kelas normal & RPM $\le$ 4000, pedal gas $\le$ 70\%, kecepatan $\le$ 120 km/jam & Tidak ada; hanya kecepatan $\ge$ 15 km/jam untuk kedua kelas \\
    \hline
    Data & Sekitar 5.000 baris per kelas & 4.536 baris normal, 4.670 baris ugal-ugalan; satu sesi per kelas \\
    \hline
    Pengemudi & 1 (peneliti) & 1 (peneliti) \\
    \hline
    Data uji & 249 jendela & 213 jendela \\
    \hline
  \end{tabular}
\end{table}

\subsection{Percobaan 1}
Tabel~\ref{tab:cm-1} dan Tabel~\ref{tab:metrik-1} menunjukkan hasil percobaan pertama pada data uji. Model mencapai akurasi 88\%. Seluruh jendela normal dikenali dengan benar, tetapi 29 dari 144 jendela ugal-ugalan (20\%) diprediksi normal. Artinya, satu dari lima kejadian ugal-ugalan lolos tanpa terdeteksi.

\begin{table}[H]
  \centering
  \caption{\textit{Confusion matrix} percobaan 1}
  \label{tab:cm-1}
  \begin{tabular}{|l|c|c|}
    \hline
     & \textbf{Prediksi ugal-ugalan} & \textbf{Prediksi normal} \\
    \hline
    \textbf{Aktual ugal-ugalan} & 115 & 29 \\
    \hline
    \textbf{Aktual normal} & 0 & 105 \\
    \hline
  \end{tabular}
\end{table}

\begin{table}[H]
  \centering
  \caption{Metrik klasifikasi percobaan 1 pada data uji}
  \label{tab:metrik-1}
  \begin{tabular}{|l|c|c|c|c|}
    \hline
    \textbf{Kelas} & \textbf{Presisi} & \textbf{\textit{Recall}} & \textbf{\textit{F1-score}} & \textbf{Jumlah jendela} \\
    \hline
    Ugal-ugalan & 1,00 & 0,80 & 0,89 & 144 \\
    \hline
    Normal & 0,78 & 1,00 & 0,88 & 105 \\
    \hline
    \multicolumn{5}{|l|}{Akurasi: 0,88} \\
    \hline
  \end{tabular}
\end{table}

Kurva pelatihan memperlihatkan akurasi latih mendekati 99,7\% sementara \textit{loss} validasi naik kembali setelah sekitar \textit{epoch} ke-8, yang menandakan \textit{overfitting}. AUC validasi terbaik sudah tercapai pada \textit{epoch} pertama, sehingga bobot yang dipakai untuk pengujian berasal dari model yang baru dilatih satu \textit{epoch}. Konvergensi secepat ini sejalan dengan filter kelas normal. Karena baris normal dengan RPM, pedal gas, atau kecepatan tinggi dibuang, kelas normal secara konstruksi tidak pernah memiliki nilai tersebut, dan model cukup mengenali batas filter. Pemeriksaan berikutnya juga menemukan bahwa struktur paket MotionSim belum terbaca dengan benar, sehingga kolom laju yaw sebenarnya berisi koordinat posisi vertikal kendaraan.

\subsection{Percobaan 2}
Berdasarkan temuan percobaan 1, dilakukan tiga perbaikan. Pertama, filter khusus kelas normal dihapus sehingga label hanya ditentukan oleh sesi rekaman. Kedua, pembacaan paket MotionSim diperbaiki sesuai struktur pada dokumentasi BeamNG dan diverifikasi dengan pengujian di dalam simulator. Ketiga, tiga fitur sudut kemudi ditambahkan. Tabel~\ref{tab:cm-2} dan Tabel~\ref{tab:metrik-2} menunjukkan hasilnya.

\begin{table}[H]
  \centering
  \caption{\textit{Confusion matrix} percobaan 2}
  \label{tab:cm-2}
  \begin{tabular}{|l|c|c|}
    \hline
     & \textbf{Prediksi ugal-ugalan} & \textbf{Prediksi normal} \\
    \hline
    \textbf{Aktual ugal-ugalan} & 105 & 0 \\
    \hline
    \textbf{Aktual normal} & 6 & 102 \\
    \hline
  \end{tabular}
\end{table}

\begin{table}[H]
  \centering
  \caption{Metrik klasifikasi percobaan 2 pada data uji}
  \label{tab:metrik-2}
  \begin{tabular}{|l|c|c|c|c|}
    \hline
    \textbf{Kelas} & \textbf{Presisi} & \textbf{\textit{Recall}} & \textbf{\textit{F1-score}} & \textbf{Jumlah jendela} \\
    \hline
    Ugal-ugalan & 0,95 & 1,00 & 0,97 & 105 \\
    \hline
    Normal & 1,00 & 0,94 & 0,97 & 108 \\
    \hline
    \multicolumn{5}{|l|}{Akurasi: 0,97} \\
    \hline
  \end{tabular}
\end{table}

Akurasi naik dari 88\% menjadi 97\%, dan \textit{F1-score} kedua kelas naik menjadi 0,97. Pola kesalahannya juga berbalik. Tidak ada jendela ugal-ugalan yang lolos (\textit{recall} 1,00), dan kesalahan yang tersisa berupa 6 jendela normal yang ditandai ugal-ugalan. Untuk sistem pemantau keselamatan, alarm palsu semacam ini lebih dapat diterima daripada perilaku berbahaya yang tidak terdeteksi. Kurva pelatihan tidak lagi menunjukkan \textit{loss} validasi yang naik, dan pelatihan berlangsung sekitar 37 \textit{epoch}.

Untuk melihat fitur yang paling menentukan, dihitung AUC setiap fitur tunggal pada dataset percobaan 2. Tabel~\ref{tab:auc-fitur} menunjukkan hasilnya. Untuk fitur yang bertanda, yaitu akselerasi, laju yaw, dan fitur kemudi, AUC dihitung dari nilai mutlaknya atau dari simpangan bakunya dalam satu jendela, karena belok ke kiri dan ke kanan sama-sama dapat menandakan perilaku ugal-ugalan.

\begin{table}[H]
  \centering
  \caption{AUC fitur tunggal pada dataset percobaan 2}
  \label{tab:auc-fitur}
  \small
  \begin{tabular}{|l|l|c|}
    \hline
    \textbf{Fitur} & \textbf{Bentuk} & \textbf{AUC} \\
    \hline
    Putaran mesin & nilai per cuplikan & 0,93 \\
    \hline
    Posisi gigi & nilai per cuplikan & 0,80 \\
    \hline
    Akselerasi longitudinal & nilai mutlak per cuplikan & 0,78 \\
    \hline
    Akselerasi lateral & simpangan baku per jendela & 0,76 \\
    \hline
    Laju yaw & simpangan baku per jendela & 0,74 \\
    \hline
    Kecepatan & nilai per cuplikan & 0,70 \\
    \hline
    Kedalaman pedal gas & nilai per cuplikan & 0,70 \\
    \hline
    Kecepatan sudut kemudi & simpangan baku per jendela & 0,68 \\
    \hline
    Sudut kemudi & simpangan baku per jendela & 0,67 \\
    \hline
  \end{tabular}
\end{table}

Putaran mesin dan posisi gigi menjadi pembeda terkuat. Pada sesi normal, 54\% cuplikan berada pada gigi 7, sedangkan pada sesi ugal-ugalan sebagian besar cuplikan berada pada gigi 3 sampai 5 dengan putaran mesin yang lebih tinggi. Hal ini sesuai dengan indikator perilaku agresif pada BAB II. Fitur gerak lateral dan kemudi tidak membedakan kelas melalui nilai sesaatnya, tetapi melalui variasinya dalam satu jendela, yang merupakan jenis pola yang ditangkap oleh lapisan konvolusi dan LSTM.

Model final memakai empat fitur inti pada Tabel~\ref{tab:fitur}, karena mobil uji hanya menyediakan data melalui OBD-II. Tiga fitur inti, yaitu putaran mesin (AUC 0,93), kecepatan (0,70), dan kedalaman pedal gas (0,70), sudah membedakan kelas pada dataset percobaan 2. Nilai AUC akselerasi longitudinal sebesar 0,78 berasal dari MotionSim, sehingga daya pisah akselerasi hasil selisih kecepatan perlu diukur ulang. Posisi gigi dengan AUC 0,80 tidak termasuk fitur inti dan hanya ditambahkan jika mobil uji menyediakannya.

\subsection{Keterbatasan Hasil Pendahuluan}
Hasil percobaan 2 belum dapat dijadikan ukuran kemampuan generalisasi model, karena tiga alasan berikut.
\begin{enumerate}
  \item Setiap kelas baru memiliki satu sesi rekaman, masing-masing sekitar 7,5 menit. Data latih, validasi, dan uji berasal dari blok-blok sesi yang sama, sehingga kondisinya mirip. Akurasi dan AUC validasi mencapai 1,00 sejak sekitar \textit{epoch} ke-6, yang menunjukkan bahwa data validasi terlalu mudah untuk memilih model terbaik.
  \item Dominasi putaran mesin dan posisi gigi perlu dipastikan berasal dari perilaku pengemudi, bukan dari perbedaan mode transmisi atau kendaraan antar sesi.
  \item Data ugal-ugalan memuat satu lonjakan akselerasi longitudinal sebesar $-18{,}9$~m/s$^{2}$ yang kemungkinan besar berasal dari tabrakan. Nilai ini sudah terpotong oleh batas $\pm 15$~m/s$^{2}$, tetapi kejadian tabrakan seharusnya tidak masuk ke dataset.
\end{enumerate}

Langkah berikutnya adalah merekam tiga sampai lima sesi per kelas dari peneliti dan dua sampai tiga pengemudi lain, dengan rute yang berbeda serta kendaraan dan mode transmisi yang sama. Setelah itu, satu sesi utuh per kelas disisihkan sebagai data uji. Sebelum pelatihan ulang, peneliti memeriksa PID yang didukung mobil uji dan mengukur laju cuplik ELM327. Model kemudian dilatih ulang dengan fitur inti pada laju cuplik tersebut memakai beberapa \textit{seed}, diuji secara \textit{real-time} di BeamNG, lalu diuji pada mobil di area tertutup.

  \cleardoublepage

  % !TEX root = ../main.tex
\chapter{JADWAL PENELITIAN}

% \G mengisi sel minggu kegiatan, \w mengosongkan sel
\newcommand{\w}{}
\newcommand{\G}{\cellcolor{gray}}
\begin{table}[H]
  \centering
  \caption{Jadwal pelaksanaan penelitian}
  \label{tbl:timeline}
  \small
  \setlength{\tabcolsep}{3pt}
  \begin{tabular}{|p{4.2cm}|c|c|c|c|c|c|c|c|c|c|c|c|c|c|c|c|}
    \hline
    \multirow{2}{*}{Kegiatan} & \multicolumn{16}{c|}{Minggu} \\
    \cline{2-17}
     & 1 & 2 & 3 & 4 & 5 & 6 & 7 & 8 & 9 & 10 & 11 & 12 & 13 & 14 & 15 & 16 \\
    \hline
    Uji coba ELM327 dan pemeriksaan PID mobil
     & \G & \G & \w & \w & \w & \w & \w & \w & \w & \w & \w & \w & \w & \w & \w & \w \\
    \hline
    Perekaman sesi simulator dari beberapa pengemudi
     & \G & \G & \G & \G & \w & \w & \w & \w & \w & \w & \w & \w & \w & \w & \w & \w \\
    \hline
    \textit{Preprocessing} dan penyesuaian fitur
     & \w & \w & \G & \G & \G & \w & \w & \w & \w & \w & \w & \w & \w & \w & \w & \w \\
    \hline
    \textit{Training} dan optimisasi model
     & \w & \w & \w & \w & \G & \G & \G & \G & \w & \w & \w & \w & \w & \w & \w & \w \\
    \hline
    Ablasi dan analisis fitur
     & \w & \w & \w & \w & \w & \w & \w & \G & \G & \w & \w & \w & \w & \w & \w & \w \\
    \hline
    Pengujian \textit{real-time} pada BeamNG
     & \w & \w & \w & \w & \w & \w & \w & \w & \G & \G & \w & \w & \w & \w & \w & \w \\
    \hline
    Pengujian pada mobil melalui OBD-II
     & \w & \w & \w & \w & \w & \w & \w & \w & \w & \w & \G & \G & \G & \w & \w & \w \\
    \hline
    Evaluasi dan perbaikan sistem
     & \w & \w & \w & \w & \w & \w & \w & \w & \w & \w & \w & \G & \G & \G & \w & \w \\
    \hline
    Penulisan buku Tugas Akhir
     & \w & \w & \w & \w & \w & \w & \w & \w & \G & \G & \G & \G & \G & \G & \G & \G \\
    \hline
  \end{tabular}
\end{table}

Tabel~\ref{tbl:timeline} menunjukkan rencana pelaksanaan penelitian selama 16 minggu. Dua minggu pertama dipakai untuk mencoba adaptor ELM327 pada mobil uji, memeriksa PID yang didukung, dan mengukur laju cuplik. Perekaman sesi simulator dari peneliti dan dua sampai tiga pengemudi lain berjalan pada empat minggu pertama. \textit{Preprocessing} dimulai pada minggu ketiga, setelah laju cuplik mobil diketahui, agar data simulator dapat langsung disesuaikan. \textit{Training} dan optimisasi model berlangsung pada minggu kelima sampai kedelapan, lalu dilanjutkan dengan ablasi dan analisis fitur. Pengujian \textit{real-time} di BeamNG dijalankan pada minggu kesembilan dan kesepuluh, dan pengujian pada mobil di area tertutup dijalankan pada minggu kesebelas sampai ketiga belas. Jika model belum memenuhi kriteria pada Tabel~\ref{tab:kriteria}, minggu kedua belas sampai keempat belas dipakai untuk memperbaiki sistem sesuai alur pada Gambar~\ref{fig:urutan}. Penulisan buku Tugas Akhir dimulai pada minggu kesembilan dan berjalan bersamaan dengan tahap pengujian.

  \cleardoublepage

  % Daftar pustaka
  \chapter*{DAFTAR PUSTAKA}
  \addcontentsline{toc}{chapter}{DAFTAR PUSTAKA}
  \renewcommand\refname{}
  \vspace{2ex}
  \renewcommand{\bibname}{}
  % Sitasi hanya ditulis di Bab 1; rujukan yang hanya disebut di Bab 2 dan 3 tetap dicetak
  \nocite{Handayani2017, Rajamani2012, BeamNG, Schuster1997, Ioffe2015, Ba2016,
    Kingma2015, Srivastava2014, Fawcett2006, Carmona2015, Singh2025, Fadzil2024,
    Sagbas2026, Ceven2022, SAEJ1979}
  \begingroup
    \def\chapter*#1{}
    \printbibliography
  \endgroup

\end{document}